\documentclass[aspectratio=169]{beamer}
\usepackage[utf8]{inputenc}
\usepackage[T1]{fontenc}
\usepackage[spanish]{babel}
\usepackage{csquotes}
\usetheme[showlogo,nosectionpage]{UU}

\usepackage{graphicx}
\usepackage{booktabs}
\usepackage{colortbl}
\usepackage{amsmath, amsfonts}
\usepackage{hyperref}
\usepackage{tikz}
\usepackage{listings}
\usepackage{fontawesome5}
\usepackage{multicol}
\usetikzlibrary{positioning, arrows.meta, shapes.geometric, fit, calc, decorations.pathreplacing}

\graphicspath{{./}{images/}}

\definecolor{dockerblue}{HTML}{2496ED}
\definecolor{terminalblack}{HTML}{1E1E1E}
\definecolor{codebg}{HTML}{F4F4F4}
\definecolor{linkedinblue}{HTML}{0A66C2}

\lstdefinelanguage{Dockerfile}{
  keywords={FROM,RUN,COPY,WORKDIR,CMD,EXPOSE,ENV,ENTRYPOINT,ADD,ARG,VOLUME,USER,LABEL,HEALTHCHECK,STOPSIGNAL,SHELL,ONBUILD},
  sensitive=true,
  comment=[l]{\#},
  morestring=[b]",
  morestring=[b]'
}
\lstdefinestyle{terminal}{
  basicstyle=\ttfamily\footnotesize,
  backgroundcolor=\color{codebg},
  keywordstyle=\bfseries\color{OTIred},
  stringstyle=\color{UUgreen},
  commentstyle=\itshape\color{UUdarkgray},
  numbers=none,
  frame=single,
  rulecolor=\color{UUmidgray},
  xleftmargin=0.5em,
  columns=fullflexible,
  keepspaces=true,
  breaklines=true
}
\lstset{
  literate={á}{{\'a}}1 {é}{{\'e}}1 {í}{{\'i}}1 {ó}{{\'o}}1 {ú}{{\'u}}1
           {Á}{{\'A}}1 {É}{{\'E}}1 {Í}{{\'I}}1 {Ó}{{\'O}}1 {Ú}{{\'U}}1
           {ñ}{{\~n}}1 {Ñ}{{\~N}}1 {ü}{{\"u}}1 {Ü}{{\"U}}1
}

\newcommand{\MainHeaderSlide}[2]{%
  \begingroup
  \setbeamercolor{background canvas}{bg=white}
  \begin{frame}[plain]
    \vfill
    \hspace*{0.9cm}%
    \begin{minipage}{0.78\paperwidth}
      {\color{UUblack}\rule{4cm}{1.3pt}}\\[14pt]
      {\raggedright\color{UUblack}\fontsize{36}{42}\selectfont\rmfamily\bfseries #1\par}
    \end{minipage}
    \vfill
  \end{frame}
  \endgroup
}

\title[Docker desde Cero]{Docker desde Cero: Crea y Despliega Aplicaciones}
\subtitle{Programa Completo — PIT 2026}
\author[Anthony Gonzalo Quispe Cordova]{Instructor: \href{https://www.linkedin.com/in/anthony-gonzalo-quispe-cordova-a39274196}{\texorpdfstring{\color{linkedinblue}\faLinkedin\ Anthony Gonzalo Quispe Cordova}{Anthony Gonzalo Quispe Cordova}}}
\institute[OTI]{Oficina de Tecnologías de la Información (OTI)}
\date{9na Edición}

\begin{document}

% ============================================================
% PORTADA
% ============================================================
\begin{frame}[plain]
  \titlepage
  \begin{tikzpicture}[remember picture, overlay]
    \node[anchor=center]
      at ([xshift=-2.05cm,yshift=-0.25cm]current page.east)
      {{\color{dockerblue}\fontsize{78}{84}\selectfont\faDocker}};
  \end{tikzpicture}
\end{frame}

% ============================================================
% ÍNDICE GENERAL
% ============================================================
\begin{frame}{Índice General del Programa}
  \tiny
  \setlength{\columnsep}{0.25cm}
  \begin{multicols}{3}
    \tableofcontents[hideallsubsections]
  \end{multicols}
\end{frame}

% ============================================================
% SESIÓN 1: CONTENEDORES DESDE CERO
% ============================================================

\begin{frame}{Objetivo de la sesión 1}
  \footnotesize
  \begin{columns}[c]
    \column{0.58\textwidth}
    \begin{block}{Al terminar la sesión 1 podrás}
      \begin{itemize}
        \item Explicar el concepto de contenerización.
        \item Explicar qué problema resuelve Docker.
        \item Ubicar Docker frente a VM e hipervisores.
        \item Identificar el flujo para instalar una VM base.
        \item Saber dónde instalar Docker y Compose desde la documentación oficial.
        \item Diferenciar imagen y contenedor.
        \item Ejecutar comandos básicos: \texttt{run}, \texttt{ps}, \texttt{stop}, \texttt{rm}.
        \item Construir y ejecutar una primera app en contenedor.
      \end{itemize}
    \end{block}
    \column{0.36\textwidth}
    \centering
    {\color{dockerblue}\fontsize{44}{50}\selectfont\faDocker}\\[0.15cm]
    {\small\textbf{Menos teoría, más terminal}}
  \end{columns}
\end{frame}

\MainHeaderSlide{Introducción y fundamentos}{Por qué Docker existe y cuándo ayuda}
\section{S1 · Introducción y Fundamentos}

\begin{frame}{El problema: «funciona en mi máquina»}
  \begin{columns}[c]
    \column{0.58\textwidth}
    \small
    \begin{itemize}
      \item La app funciona en la laptop del desarrollador.
      \item En otro equipo falla por versiones distintas.
      \item En producción aparecen librerías faltantes o configuraciones diferentes.
      \item El equipo pierde tiempo reconstruyendo el entorno.
    \end{itemize}
    \begin{alertblock}{Clave}
      Docker busca que el entorno viaje junto con la aplicación.
    \end{alertblock}
    \column{0.35\textwidth}
    \centering
    {\color{OTIred}\fontsize{38}{44}\selectfont\faExclamationTriangle}\\[0.1cm]
    {\large\textbf{«En mi PC funciona»}}
  \end{columns}
\end{frame}

\begin{frame}{Solución: la contenerización}
  \begin{columns}[c]
    \column{0.56\textwidth}
    \small
    \begin{itemize}
      \item Aislar la aplicación en un \textbf{contenedor} reproducible.
      \item Empaquetar código, dependencias y configuración en una \textbf{imagen}.
      \item Compartir esa imagen con el equipo o con un registro central.
      \item Ejecutarla en distintos entornos con menos diferencias.
    \end{itemize}
    \begin{block}{Resultado}
      La app deja de depender de «lo que justo estaba instalado» en una máquina.
    \end{block}
    \column{0.38\textwidth}
    \centering
    {\color{dockerblue}\fontsize{44}{50}\selectfont\faDocker}\\[0.25cm]
    {\large\textbf{App + dependencias}}
  \end{columns}
\end{frame}

\begin{frame}{Qué es la contenerización}
  \begin{columns}[T]
    \column{0.54\textwidth}
    \small
    \begin{block}{En palabras simples}
      La contenerización empaqueta una aplicación con lo necesario para ejecutarla de forma uniforme en diferentes entornos.
    \end{block}
    \begin{itemize}
      \item Aplicación + dependencias.
      \item Configuración mínima reproducible.
      \item Ejecución aislada del resto del sistema.
      \item Menos problemas de compatibilidad.
    \end{itemize}
    \column{0.40\textwidth}
    \centering
    \begin{tikzpicture}[
      box/.style={draw, rounded corners=6pt, minimum width=3.2cm, minimum height=1.2cm, align=center, text width=3.0cm, font=\small},
      arr/.style={-{Stealth[length=3mm]}, thick, OTIred}
    ]
      \node[box, fill=orange!15, draw=orange!70] (app) at (0,0) {\faCode\ \textbf{App}};
      \node[box, fill=blue!10, draw=blue!60, below=0.55cm of app] (deps) {\faBox\ \textbf{Dependencias}};
      \node[box, fill=green!12, draw=green!60!black, below=0.55cm of deps] (container) {\faArchive\ \textbf{Contenedor}};
      \draw[arr] (app) -- (deps);
      \draw[arr] (deps) -- (container);
    \end{tikzpicture}
  \end{columns}
\end{frame}

\begin{frame}{La pregunta natural}
  \begin{columns}[c]
    \column{0.48\textwidth}
    \small
    \begin{block}{Ya tenemos la idea}
      Contenerizar es empaquetar una app con sus dependencias para ejecutarla igual en distintos entornos.
    \end{block}
    \vspace{0.18cm}
    \begin{alertblock}{Ahora falta la herramienta}
      \textbf{Docker} es una de las formas más usadas para construir, compartir y ejecutar esos contenedores.
    \end{alertblock}
    \vspace{0.1cm}
    {\color{dockerblue}\bfseries Entonces... ¿qué es Docker?}
    \column{0.46\textwidth}
    \centering
    {\color{dockerblue}\fontsize{64}{70}\selectfont\faDocker}\\[0.2cm]
    {\large\bfseries Docker =}\\[0.1cm]
    {\small Construir · Compartir · Ejecutar}\\[0.1cm]
    {\small\itshape contenedores}
  \end{columns}
\end{frame}

\begin{frame}{Entonces... ¿qué es Docker?}
  \begin{center}
    {\color{UUblack}\fontsize{25}{30}\selectfont\rmfamily\bfseries Si una VM virtualiza una máquina...}\\[0.30cm]
    {\color{dockerblue}\fontsize{30}{35}\selectfont\rmfamily\bfseries Docker ejecuta aplicaciones}\\[0.08cm]
    {\color{dockerblue}\fontsize{30}{35}\selectfont\rmfamily\bfseries en contenedores}\\[0.24cm]
    {\color{OTIred}\rule{4cm}{1.3pt}}
  \end{center}
\end{frame}

\begin{frame}{Docker en una frase}
  \begin{center}
  \begin{tikzpicture}[
    box/.style={draw, rounded corners=6pt, minimum width=3.5cm, minimum height=1.35cm, align=center, text width=3.3cm, font=\small},
    arr/.style={-{Stealth[length=3mm]}, thick, OTIred}
  ]
    \node[box, fill=OTIred!10, draw=OTIred!70] (code) at (0,0) {\faCode\ \textbf{Código}\\App Flask};
    \node[box, fill=orange!15, draw=orange!70, right=1.0cm of code] (deps) {\faBox\ \textbf{Dependencias}\\Python + librerías};
    \node[box, fill=blue!10, draw=blue!60, right=1.0cm of deps] (img) {\faDocker\ \textbf{Imagen}\\Paquete reproducible};
    \draw[arr] (code) -- (deps);
    \draw[arr] (deps) -- (img);
  \end{tikzpicture}
  \end{center}
  \vspace{0.2cm}
  \begin{block}{Definición práctica}
    Docker permite empaquetar una aplicación con sus dependencias para ejecutarla de forma consistente en diferentes equipos.
  \end{block}
\end{frame}

\begin{frame}{Docker: introducción práctica}
  \begin{columns}[c]
    \column{0.58\textwidth}
    \footnotesize
    \begin{itemize}
      \item Docker es una plataforma \textit{open source} para crear, compartir y ejecutar contenedores.
      \item Implementa la idea de contenerización en el flujo diario de desarrollo.
      \item Facilita pruebas, despliegues y trabajo en equipo.
      \item Se apoya en imágenes, contenedores, registros y la CLI de Docker.
    \end{itemize}
    \begin{alertblock}{Idea central}
      Docker no reemplaza toda VM: resuelve la ejecución reproducible de aplicaciones.
    \end{alertblock}
    \column{0.36\textwidth}
    \centering
    {\color{dockerblue}\fontsize{58}{64}\selectfont\faDocker}\\[0.25cm]
    {\Large\textbf{Docker}}
  \end{columns}
\end{frame}

\begin{frame}{Qué sí veremos y qué no veremos hoy}
  \begin{columns}[T]
    \column{0.48\textwidth}
    \begin{exampleblock}{Hoy sí}
      \begin{itemize}
        \item Concepto de Docker.
        \item Concepto de contenerización.
        \item Ruta oficial de instalación.
        \item Imagen vs contenedor.
        \item Comandos esenciales.
        \item Primera app con Dockerfile.
      \end{itemize}
    \end{exampleblock}
    \column{0.48\textwidth}
    \begin{alertblock}{Hoy no}
      \begin{itemize}
        \item Kubernetes, OCI profundo o cgroups.
        \item YAML avanzado.
        \item Redes y volúmenes a detalle.
        \item Producción con Nginx.
      \end{itemize}
    \end{alertblock}
  \end{columns}
\end{frame}

\begin{frame}{Pero... ¿cómo funciona?}
  \begin{center}
    {\color{UUblack}\fontsize{19}{24}\selectfont\rmfamily\bfseries Entendido... ¿Pero cómo funciona?}\\[0.35cm]
    {\color{UUdarkgray}\fontsize{18}{23}\selectfont\itshape ¿Es Docker simplemente otra...}\\[0.28cm]
    {\color{dockerblue}\fontsize{32}{38}\selectfont\rmfamily\bfseries Máquina Virtual?}\\[0.28cm]
    {\color{OTIred}\rule{4.2cm}{1.3pt}}
  \end{center}
\end{frame}

\MainHeaderSlide{Virtualización\\e hipervisores}{Cómo se diferencia Docker de una máquina virtual}
\section{S1 · Virtualización e Hipervisores}

\begin{frame}{Antes de Docker: máquinas virtuales}
  \begin{columns}[T]
    \column{0.56\textwidth}
    \footnotesize
    \begin{itemize}
      \item Una VM ejecuta un sistema operativo completo sobre hardware virtual.
      \item Permite aislar entornos: Linux, Windows Server, laboratorios, pruebas.
      \item Es útil cuando necesitas simular una máquina completa.
      \item Consume más recursos que un contenedor porque incluye su propio kernel.
    \end{itemize}
    \begin{alertblock}{Idea clave}
      La VM virtualiza una máquina; Docker empaqueta procesos aislados.
    \end{alertblock}
    \column{0.38\textwidth}
    \centering
    {\color{UUpurple}\fontsize{42}{48}\selectfont\faServer}\\[0.2cm]
    {\normalsize\textbf{Máquina completa}}
  \end{columns}
\end{frame}

\begin{frame}{Qué es un hipervisor}
  \begin{center}
  \begin{tikzpicture}[
    box/.style={draw, rounded corners=6pt, minimum width=4.7cm, minimum height=1.4cm, align=center, text width=4.4cm, font=\small},
  ]
    \node[box, fill=OTIred!10, draw=OTIred!70] (type1) at (0,0) {
      {\color{OTIred}\faMicrochip\ \textbf{Tipo 1}}\\[3pt]
      Corre directo sobre el hardware. Ejemplos: VMware ESXi, Hyper-V Server, Proxmox.
    };
    \node[box, fill=blue!10, draw=blue!60, right=1.0cm of type1] (type2) {
      {\color{dockerblue}\faLaptop\ \textbf{Tipo 2}}\\[3pt]
      Corre como aplicación del sistema host. Ejemplos: VirtualBox, VMware Workstation.
    };
  \end{tikzpicture}
  \end{center}
  \vspace{0.15cm}
  \begin{block}{Para este curso}
    Si necesitas un laboratorio local, una VM con Linux puede servir como entorno limpio para instalar Docker y practicar sin afectar tu sistema principal.
  \end{block}
\end{frame}

\begin{frame}{La siguiente duda}
  \begin{center}
    {\color{UUblack}\fontsize{26}{32}\selectfont\rmfamily\bfseries ¿Y }%
    {\color{OTIred}\fontsize{26}{32}\selectfont\rmfamily\bfseries Docker}%
    {\color{UUblack}\fontsize{26}{32}\selectfont\rmfamily\bfseries ?}\\[0.38cm]
    {\color{UUdarkgray}\fontsize{20}{25}\selectfont\itshape ¿Es un hipervisor?}\\[0.30cm]
    {\color{UUdarkgray}\fontsize{19}{24}\selectfont\itshape ¿De qué tipo?}\\[0.30cm]
    {\color{OTIred}\rule{4.0cm}{1.3pt}}
  \end{center}
\end{frame}

\begin{frame}{Docker vs máquina virtual}
  \begin{center}
    \begin{tikzpicture}[scale=0.64, transform shape,
      box/.style={draw, rounded corners=6pt, minimum width=2.8cm, minimum height=1.2cm, align=center, font=\scriptsize},
    ]
      % VM side
      \node[font=\small\bfseries] at (-1.5, 6.5) {\color{OTIred} Máquina Virtual};
      \node[box, fill=OTIred!10, draw=OTIred!70] at (-1.5, 5.5) {App};
      \node[box, fill=OTIred!10, draw=OTIred!70] at (-1.5, 4.2) {Bins / Libs};
      \node[box, fill=OTIred!20, draw=OTIred!70] at (-1.5, 2.9) {Guest OS};
      \node[box, fill=UUdarkgray!20, draw=UUdarkgray] at (-1.5, 1.6) {Hypervisor};
      \node[box, fill=UUlightgray, draw=UUdarkgray] at (-1.5, 0.3) {Infrastructure};
      
      % Container side
      \node[font=\small\bfseries] at (1.5, 6.5) {\color{dockerblue} Contenedor};
      \node[box, fill=blue!10, draw=blue!60] at (1.5, 5.5) {App};
      \node[box, fill=blue!10, draw=blue!60] at (1.5, 4.2) {Bins / Libs};
      \node[box, fill=green!12, draw=green!60!black] at (1.5, 2.9) {Docker Engine};
      \node[box, fill=UUlightgray, draw=UUdarkgray] at (1.5, 0.3) {Host OS + Infrastructure};
      
      % Braces
      \draw[OTIred, thick, decorate, decoration={brace, amplitude=5pt}] (-3.2, 4.7) -- (-3.2, 5.8) node[midway, left, font=\tiny] {App};
      \draw[dockerblue, thick, decorate, decoration={brace, amplitude=5pt}] (3.2, 4.7) -- (3.2, 5.8) node[midway, right, font=\tiny] {App};
    \end{tikzpicture}
  \end{center}
  \vspace{-0.15cm}
  \begin{alertblock}{Diferencia clave}
    \small
    La VM incluye un sistema operativo completo (Guest OS). El contenedor comparte el kernel del host y solo incluye lo que la app necesita.
  \end{alertblock}
\end{frame}

\begin{frame}{VM vs contenedor: decisión rápida}
  \scriptsize
  \begin{tabular}{p{0.25\textwidth} p{0.31\textwidth} p{0.33\textwidth}}
    \toprule
    Criterio & Máquina virtual & Contenedor \\
    \midrule
    Sistema operativo & Incluye SO completo & Comparte kernel del host \\
    Arranque & Más lento & Rápido \\
    Consumo & Mayor RAM y disco & Menor consumo \\
    Aislamiento & Fuerte a nivel máquina & Aislado a nivel proceso \\
    Uso común & Servidores y laboratorios & Apps reproducibles \\
    \bottomrule
  \end{tabular}
\end{frame}

\begin{frame}{Instalando una VM: flujo recomendado}
  \begin{center}
  \begin{tikzpicture}[
    stage/.style={draw, rounded corners=6pt, minimum width=2.35cm, minimum height=1.2cm, align=center, text width=2.15cm, font=\scriptsize, fill=UUlightgray, draw=UUdarkgray},
    arr/.style={-{Stealth[length=3mm]}, thick, OTIred}
  ]
    \node[stage] (iso) {\faDownload\\\textbf{ISO}\\Ubuntu Server};
    \node[stage, right=0.45cm of iso] (vm) {\faCog\\\textbf{Crear VM}\\CPU, RAM, disco};
    \node[stage, right=0.45cm of vm] (net) {\faNetworkWired\\\textbf{Red}\\NAT o bridge};
    \node[stage, right=0.45cm of net] (install) {\faPlay\\\textbf{Instalar}\\Usuario y SSH};
    \node[stage, right=0.45cm of install] (docker) {\faDocker\\\textbf{Docker}\\Practicar};
    \draw[arr] (iso) -- (vm);
    \draw[arr] (vm) -- (net);
    \draw[arr] (net) -- (install);
    \draw[arr] (install) -- (docker);
  \end{tikzpicture}
  \end{center}
  \begin{alertblock}{Configuración sugerida}
    2 CPU, 4 GB RAM, 25 GB disco, red NAT para empezar. Si necesitas acceso desde otros equipos, usa bridge.
  \end{alertblock}
\end{frame}

\MainHeaderSlide{Instalación de Docker}{Preparar el entorno usando solo documentación oficial}
\section{S1 · Instalación de Docker}

\begin{frame}{Qué debes instalar}
  \begin{center}
  \begin{tikzpicture}[
    box/.style={draw, rounded corners=6pt, minimum width=3.5cm, minimum height=1.35cm, align=center, text width=3.25cm, font=\small},
    arr/.style={-{Stealth[length=3mm]}, thick, OTIred}
  ]
    \node[box, fill=blue!10, draw=blue!60] (docker) {\faDocker\ \textbf{Docker}\\Engine o Desktop};
    \node[box, fill=green!12, draw=green!60!black, right=0.8cm of docker] (compose) {\faLayerGroup\ \textbf{Docker Compose}\\Orquestación local};
    \node[box, fill=orange!15, draw=orange!70, right=0.8cm of compose] (verify) {\faCheckCircle\ \textbf{Validación}\\Entorno listo};
    \draw[arr] (docker) -- (compose);
    \draw[arr] (compose) -- (verify);
  \end{tikzpicture}
  \end{center}
  \vspace{0.2cm}
  \begin{block}{Criterio del curso}
    No memorices comandos de instalación: pueden cambiar por sistema operativo. Usa siempre la guía oficial y valida según tu plataforma.
  \end{block}
\end{frame}

\begin{frame}{Ruta oficial según tu sistema}
  \small
  \begin{columns}[T]
    \column{0.48\textwidth}
    \begin{exampleblock}{Windows y macOS}
      \begin{itemize}
        \item Instalar \textbf{Docker Desktop}.
        \item Incluye Docker Engine, CLI y Docker Compose.
        \item Revisar requisitos de WSL 2 en Windows.
      \end{itemize}
      \vspace{0.1cm}
      \href{https://docs.docker.com/desktop/}{\faExternalLink*\ docs.docker.com/desktop}
    \end{exampleblock}
    \column{0.48\textwidth}
    \begin{block}{Linux o VM Linux}
      \begin{itemize}
        \item Instalar \textbf{Docker Engine} según distribución.
        \item Revisar pasos post-instalación.
        \item Instalar Compose desde la guía oficial.
      \end{itemize}
      \vspace{0.1cm}
      \href{https://docs.docker.com/engine/install/}{\faExternalLink*\ docs.docker.com/engine/install}
    \end{block}
  \end{columns}
\end{frame}

\begin{frame}{Docker Compose y enlaces oficiales}
  \footnotesize
  \begin{tabular}{p{0.29\textwidth} p{0.36\textwidth} p{0.25\textwidth}}
    \toprule
    Tema & Cuándo usarlo & Documentación \\
    \midrule
    Docker Desktop & Windows/macOS o experiencia integrada & \href{https://docs.docker.com/desktop/}{Docker Desktop} \\
    Docker Engine & Linux nativo o VM Linux & \href{https://docs.docker.com/engine/install/}{Docker Engine} \\
    Post-instalación Linux & Permisos, servicio y ajustes posteriores & \href{https://docs.docker.com/engine/install/linux-postinstall/}{Linux post-install} \\
    Docker Compose & Aplicaciones con varios servicios & \href{https://docs.docker.com/compose/install/}{Compose install} \\
    \bottomrule
  \end{tabular}
  \vspace{0.25cm}
  \begin{alertblock}{Regla práctica}
    Si una guía externa contradice la documentación oficial, se prioriza Docker Docs.
  \end{alertblock}
\end{frame}

\MainHeaderSlide{Imágenes y contenedores}{La diferencia que debes dominar antes de escribir comandos}
\section{S1 · Imágenes y Contenedores}

\begin{frame}{Imagen vs contenedor}
  \begin{center}
  \begin{tikzpicture}[
    box/.style={draw, rounded corners=6pt, minimum width=4.8cm, minimum height=2.0cm, align=center, text width=4.5cm, font=\small},
    arr/.style={-{Stealth[length=3mm]}, very thick, dockerblue}
  ]
    \node[box, fill=blue!10, draw=blue!60] (image) at (0,0) {
      {\color{dockerblue}\faArchive\ \textbf{Imagen}}\\[4pt]
      Plantilla inmutable con sistema base, dependencias y código.
    };
    \node[box, fill=green!12, draw=green!60!black, right=1.2cm of image] (container) {
      {\color{green!60!black}\faPlayCircle\ \textbf{Contenedor}}\\[4pt]
      Instancia en ejecución creada desde una imagen.
    };
    \draw[arr] (image) -- node[above, font=\scriptsize\bfseries] {docker run} (container);
  \end{tikzpicture}
  \end{center}
  \begin{alertblock}{Regla mental}
    La imagen es la receta; el contenedor es el plato servido y ejecutándose.
  \end{alertblock}
\end{frame}

\begin{frame}{Ciclo de vida básico}
  \begin{center}
  \begin{tikzpicture}[
    stage/.style={draw, rounded corners=6pt, minimum width=2.8cm, minimum height=1.25cm, align=center, fill=UUlightgray, draw=UUdarkgray},
    arr/.style={-{Stealth[length=3mm]}, thick, OTIred}
  ]
    \node[stage] (pull) {\faDownload\\\textbf{Obtener}\\\texttt{docker pull}};
    \node[stage, right=0.65cm of pull] (run) {\faPlay\\\textbf{Ejecutar}\\\texttt{docker run}};
    \node[stage, right=0.65cm of run] (ps) {\faList\\\textbf{Ver}\\\texttt{docker ps}};
    \node[stage, right=0.65cm of ps] (stop) {\faStop\\\textbf{Detener}\\\texttt{docker stop}};
    \draw[arr] (pull) -- (run);
    \draw[arr] (run) -- (ps);
    \draw[arr] (ps) -- (stop);
  \end{tikzpicture}
  \end{center}
  \vspace{0.2cm}
  \begin{block}{Lo mínimo para operar}
    Ejecutar, listar, detener, eliminar contenedores y revisar imágenes locales.
  \end{block}
\end{frame}

\MainHeaderSlide{Comandos esenciales}{Run, ps, stop, rm e images en práctica}
\section{S1 · Comandos Esenciales}

\begin{frame}[fragile]{Laboratorio 1: comprobar Docker}
  \begin{columns}[T]
    \column{0.50\textwidth}
\begin{lstlisting}[style=terminal]
docker --version
docker info
docker run hello-world
\end{lstlisting}
    \column{0.46\textwidth}
    \begin{block}{Qué validamos}
      \begin{itemize}
        \item La CLI responde.
        \item El Engine está activo.
        \item Docker puede descargar y ejecutar imágenes.
      \end{itemize}
    \end{block}
  \end{columns}
\end{frame}

\begin{frame}[fragile]{Laboratorio 2: ejecutar Nginx}
  \begin{columns}[T]
    \column{0.55\textwidth}
\begin{lstlisting}[style=terminal]
docker run --name web-demo -d -p 8080:80 nginx

docker ps

# Abrir en navegador:
# http://localhost:8080
\end{lstlisting}
    \column{0.40\textwidth}
    \begin{exampleblock}{Qué significa}
      \begin{itemize}
        \item \texttt{--name}: nombre del contenedor.
        \item \texttt{-d}: ejecuta en segundo plano.
        \item \texttt{-p 8080:80}: publica el puerto.
        \item \texttt{nginx}: imagen usada.
      \end{itemize}
    \end{exampleblock}
  \end{columns}
\end{frame}

\begin{frame}[fragile]{Laboratorio 3: detener y limpiar}
  \begin{columns}[T]
    \column{0.55\textwidth}
\begin{lstlisting}[style=terminal]
docker stop web-demo
docker ps -a
docker rm web-demo
docker images
\end{lstlisting}
    \column{0.40\textwidth}
    \begin{alertblock}{Buenas prácticas}
      \begin{itemize}
        \item No acumules contenedores detenidos.
        \item Usa nombres claros en laboratorios.
        \item Verifica antes de borrar.
      \end{itemize}
    \end{alertblock}
  \end{columns}
\end{frame}

\begin{frame}{Mapa rápido de comandos}
  \footnotesize
  \begin{tabular}{p{0.23\textwidth} p{0.32\textwidth} p{0.34\textwidth}}
    \toprule
    Comando & Para qué sirve & Ejemplo \\
    \midrule
    \texttt{docker run} & Crea y ejecuta un contenedor & \texttt{docker run nginx} \\
    \texttt{docker ps} & Lista contenedores activos & \texttt{docker ps -a} \\
    \texttt{docker stop} & Detiene un contenedor & \texttt{docker stop web-demo} \\
    \texttt{docker rm} & Elimina contenedores detenidos & \texttt{docker rm web-demo} \\
    \texttt{docker images} & Lista imágenes locales & \texttt{docker images} \\
    \bottomrule
  \end{tabular}
\end{frame}

\MainHeaderSlide{Primera aplicación}{Construir y ejecutar una app Flask mínima}
\section{S1 · Primera Aplicación}

\begin{frame}{Proyecto práctico de hoy}
  \begin{center}
  \begin{tikzpicture}[
    box/.style={draw, rounded corners=6pt, minimum width=2.85cm, minimum height=1.35cm, align=center, text width=2.65cm, font=\scriptsize},
    arr/.style={-{Stealth[length=3mm]}, thick, OTIred}
  ]
    \node[box, fill=green!12, draw=green!60!black] (flask) {\faPython\ \textbf{Flask}\\Aplicación web};
    \node[box, fill=blue!10, draw=blue!60, right=0.55cm of flask] (dockerfile) {\faFileCode\ \textbf{Dockerfile}\\Receta de imagen};
    \node[box, fill=orange!15, draw=orange!70, right=0.55cm of dockerfile] (image) {\faArchive\ \textbf{Imagen}\\\texttt{mi-flask:v1}};
    \node[box, fill=OTIred!10, draw=OTIred!70, right=0.55cm of image] (container) {\faPlay\ \textbf{Contenedor}\\Puerto 5000};
    \draw[arr] (flask) -- (dockerfile);
    \draw[arr] (dockerfile) -- (image);
    \draw[arr] (image) -- (container);
  \end{tikzpicture}
  \end{center}
  \begin{block}{Meta}
    Crear una app mínima, construir su imagen y verla funcionando desde el navegador.
  \end{block}
\end{frame}

\begin{frame}[fragile]{Paso 1: crear archivos del proyecto}
\begin{lstlisting}[style=terminal]
cd codigos/sesion1

# Archivos ya preparados para la practica:
# app.py
# requirements.txt
# Dockerfile
\end{lstlisting}
  \begin{alertblock}{Estructura esperada}
    Los archivos de la app están en \texttt{codigos/sesion1} junto al material de esta sesión.
  \end{alertblock}
\end{frame}

\begin{frame}[fragile]{Paso 2: código de la app Flask}
\begin{lstlisting}[language=Python]
from flask import Flask

app = Flask(__name__)

@app.route("/")
def home():
    return "Hola Docker desde PIT 2026"

if __name__ == "__main__":
    app.run(host="0.0.0.0", port=5000)
\end{lstlisting}
\end{frame}

\begin{frame}[fragile]{Paso 3: dependencias}
\begin{lstlisting}[style=terminal]
# requirements.txt
flask==3.0.3
\end{lstlisting}
  \begin{block}{Por qué declararlas}
    El contenedor debe poder instalar lo necesario sin depender de lo que tenga la laptop del alumno.
  \end{block}
\end{frame}

\begin{frame}[fragile]{Paso 4: Dockerfile mínimo}
\begin{lstlisting}[language=Dockerfile]
FROM python:3.12-slim

WORKDIR /app

COPY requirements.txt .
RUN pip install --no-cache-dir -r requirements.txt

COPY app.py .

EXPOSE 5000
CMD ["python", "app.py"]
\end{lstlisting}
\end{frame}

\begin{frame}[fragile]{Paso 5: construir la imagen}
\begin{lstlisting}[style=terminal]
docker build -t mi-flask:v1 .

docker images
\end{lstlisting}
  \begin{exampleblock}{Lectura del comando}
    \texttt{-t mi-flask:v1} asigna nombre y versión a la imagen. El punto final indica que el contexto de build es la carpeta actual.
  \end{exampleblock}
\end{frame}

\begin{frame}[fragile]{Paso 6: ejecutar la aplicación}
\begin{lstlisting}[style=terminal]
docker run --name flask-demo -d -p 5000:5000 mi-flask:v1

docker ps

# Abrir:
# http://localhost:5000
\end{lstlisting}
  \begin{block}{Resultado esperado}
    El navegador debe mostrar: \texttt{Hola Docker desde PIT 2026}
  \end{block}
\end{frame}

\begin{frame}[fragile]{Paso 7: limpiar el laboratorio}
\begin{lstlisting}[style=terminal]
docker stop flask-demo
docker rm flask-demo

# Opcional: borrar la imagen
docker rmi mi-flask:v1
\end{lstlisting}
  \begin{alertblock}{Regla práctica}
    Primero detienes el contenedor, luego lo eliminas. La imagen queda disponible hasta que decidas borrarla.
  \end{alertblock}
\end{frame}

\begin{frame}{Errores frecuentes en la primera práctica}
  \footnotesize
  \begin{tabular}{p{0.30\textwidth} p{0.30\textwidth} p{0.30\textwidth}}
    \toprule
    Error & Causa probable & Solución rápida \\
    \midrule
    Puerto ocupado & Otro proceso usa 5000 & Cambiar a \texttt{-p 5001:5000} \\
    No abre en navegador & Contenedor detenido & Revisar \texttt{docker ps -a} \\
    Imagen no encontrada & Nombre mal escrito & Revisar \texttt{docker images} \\
    Build falla & Archivo faltante o typo & Verificar carpeta y Dockerfile \\
    \bottomrule
  \end{tabular}
\end{frame}

\begin{frame}{Checklist de aprendizaje — Sesión 1}
  \begin{itemize}
    \item Puedo explicar qué es un hipervisor y para qué sirve una VM.
    \item Puedo describir el flujo básico para instalar una VM de laboratorio.
    \item Puedo ubicar la documentación oficial para instalar Docker y Compose.
    \item Puedo explicar qué es la contenerización y qué problema resuelve.
    \item Puedo explicar por qué Docker evita diferencias entre entornos.
    \item Puedo diferenciar imagen y contenedor con un ejemplo.
    \item Puedo ejecutar, listar, detener y eliminar contenedores.
    \item Puedo construir una imagen propia con \texttt{docker build}.
    \item Puedo ejecutar mi primera app usando \texttt{docker run -p}.
  \end{itemize}
\end{frame}

% ============================================================
% TRANSICIÓN A SESIÓN 2
% ============================================================

\begin{frame}{Resumen de la Sesión 1}
  \footnotesize
  \begin{enumerate}
    \item Una VM virtualiza una máquina completa mediante un hipervisor.
    \item La instalación de Docker y Compose se valida con la documentación oficial.
    \item La contenerización empaqueta aplicaciones y dependencias.
    \item Docker permite ejecutar entornos reproducibles.
    \item Una imagen es una plantilla; un contenedor es una ejecución de esa imagen.
    \item Ya ejecutamos comandos base y una primera aplicación Flask en contenedor.
  \end{enumerate}
  \vspace{0.08cm}
  \begin{alertblock}{Ahora: Sesión 2}
    \textbf{Dockerfile Profesional:} instrucciones clave, capas, caché y buenas prácticas.
  \end{alertblock}
\end{frame}

% ============================================================
% SESIÓN 2: DOCKERFILE PROFESIONAL
% ============================================================

\begin{frame}{Objetivo de la sesión 2}
  \footnotesize
  \begin{columns}[c]
    \column{0.58\textwidth}
    \begin{block}{Al terminar la sesión 2 podrás}
      \begin{itemize}
        \item Escribir un Dockerfile con las instrucciones clave: FROM, RUN, COPY, ENV, EXPOSE y CMD.
        \item Explicar cómo funcionan las capas y cómo afectan el rendimiento del build.
        \item Aprovechar la caché de Docker para builds más rápidos.
        \item Usar .dockerignore para evitar archivos innecesarios en la imagen.
        \item Elegir imágenes base livianas (slim, alpine) según el caso.
        \item Publicar imágenes en Docker Hub con tagging y versionado.
      \end{itemize}
    \end{block}
    \column{0.36\textwidth}
    \centering
    {\color{dockerblue}\fontsize{44}{50}\selectfont\faDocker}\\[0.15cm]
    {\small\textbf{De Dockerfile básico a imagen profesional}}
  \end{columns}
\end{frame}

\MainHeaderSlide{Recordatorio de la sesión\\anterior}{Lo que ya sabemos sobre imágenes y contenedores}
\section{S2 · Recordatorio}

\begin{frame}{Lo que ya dominamos}
  \begin{columns}[c]
    \column{0.55\textwidth}
    \footnotesize
    \begin{itemize}
      \item Contenerización: empaquetar app + dependencias.
      \item Diferencia entre imagen (plantilla) y contenedor (ejecución).
      \item Docker resuelve el problema de \textit{«en mi máquina funciona»}.
      \item Comandos esenciales: \texttt{run}, \texttt{ps}, \texttt{stop}, \texttt{rm}, \texttt{images}.
      \item Construimos y ejecutamos una app Flask en contenedor.
    \end{itemize}
    \begin{exampleblock}{Base para hoy}
      Sabemos ejecutar contenedores desde imágenes existentes. Ahora crearemos nuestras propias imágenes de forma profesional.
    \end{exampleblock}
    \column{0.38\textwidth}
    \centering
    \begin{tikzpicture}[
      box/.style={draw, rounded corners=6pt, minimum width=3cm, minimum height=1.1cm, align=center, font=\small},
      arr/.style={-{Stealth[length=3mm]}, thick, OTIred}
    ]
      \node[box, fill=blue!10, draw=blue!60] (img) {\faArchive\\\textbf{Imagen}};
      \node[box, fill=green!12, draw=green!60!black, below=0.6cm of img] (cont) {\faPlay\\\textbf{Contenedor}};
      \draw[arr] (img) -- node[right, font=\scriptsize] {docker run} (cont);
    \end{tikzpicture}
  \end{columns}
\end{frame}

\begin{frame}{¿Qué es un Dockerfile?}
  \begin{center}
    \begin{tikzpicture}[
      box/.style={draw, rounded corners=6pt, minimum width=3cm, minimum height=1.4cm, align=center, font=\small},
      arr/.style={-{Stealth[length=3mm]}, thick, dockerblue}
    ]
      \node[box, fill=orange!15, draw=orange!70] (dockerfile) at (0,0) {\faFileCode\\\textbf{Dockerfile}\\\small Receta de imagen};
      \node[box, fill=blue!10, draw=blue!60, right=1cm of dockerfile] (build) {\faCog\\\textbf{docker build}\\\small Construye la imagen};
      \node[box, fill=green!12, draw=green!60!black, right=1cm of build] (img) {\faArchive\\\textbf{Imagen}\\\small Paquete reproducible};
      \draw[arr] (dockerfile) -- (build);
      \draw[arr] (build) -- (img);
    \end{tikzpicture}
  \end{center}
  \vspace{0.2cm}
  \begin{block}{Definición}
    Un Dockerfile es un archivo de texto con instrucciones que Docker lee para construir una imagen automáticamente. Es la \textbf{receta} que define todo lo que necesita tu aplicación para ejecutarse.
  \end{block}
\end{frame}

\MainHeaderSlide{Instrucciones clave\\del Dockerfile}{FROM, RUN, COPY, ENV, EXPOSE y CMD}
\section{S2 · Instrucciones Clave del Dockerfile}

\begin{frame}[fragile]{FROM: el punto de partida}
  \begin{columns}[T]
    \column{0.54\textwidth}
\begin{lstlisting}[language=Dockerfile]
FROM python:3.12-slim
\end{lstlisting}
    \vspace{0.2cm}
    \begin{itemize}
      \item Define la imagen base sobre la que se construye.
      \item Es la \textbf{primera instrucción} de todo Dockerfile.
      \item Puede ser una imagen oficial (python, node, nginx) o una imagen propia.
    \end{itemize}
    \column{0.42\textwidth}
    \begin{exampleblock}{Consejo}
      Usa imágenes oficiales siempre que puedas. Son mantenidas, actualizadas y más seguras que imágenes de terceros.
    \end{exampleblock}
    \begin{alertblock}{Versiones}
      Evita \texttt{:latest} en producción. Fija una versión concreta (\texttt{:3.12-slim}) para builds reproducibles.
    \end{alertblock}
  \end{columns}
\end{frame}

\begin{frame}[fragile]{RUN: ejecutar comandos durante el build}
  \begin{columns}[T]
    \column{0.54\textwidth}
\begin{lstlisting}[language=Dockerfile]
RUN apt-get update && \
    apt-get install -y curl && \
    rm -rf /var/lib/apt/lists/*
\end{lstlisting}
    \vspace{0.2cm}
    \begin{itemize}
      \item Ejecuta comandos \textbf{durante la construcción} de la imagen.
      \item Se usa para instalar paquetes, crear carpetas, configurar el sistema.
      \item Cada RUN crea una \textbf{nueva capa} en la imagen.
    \end{itemize}
    \column{0.42\textwidth}
    \begin{block}{Buena práctica}
      Encadena comandos con \texttt{\&\&} para reducir capas y limpiar cachés en la misma instrucción.
    \end{block}
    \begin{alertblock}{Regla de oro}
      Limpia archivos temporales (\texttt{apt lists}, \texttt{pip cache}) en el mismo RUN para no arrastrar basura a la imagen final.
    \end{alertblock}
  \end{columns}
\end{frame}

\begin{frame}[fragile]{COPY vs ADD}
  \footnotesize
  \begin{columns}[T]
    \column{0.48\textwidth}
\begin{lstlisting}[language=Dockerfile, basicstyle=\ttfamily\scriptsize]
# COPY: copia archivos locales
COPY requirements.txt .
COPY app.py /app/
\end{lstlisting}
    \vspace{0.15cm}
    \begin{itemize}
      \item Copia archivos o carpetas desde el contexto de build (tu PC) hacia la imagen.
      \item Es la opción \textbf{recomendada} para la mayoría de los casos.
    \end{itemize}
    \column{0.48\textwidth}
\begin{lstlisting}[language=Dockerfile, basicstyle=\ttfamily\scriptsize]
# ADD: copia + descomprime + URLs
ADD app.tar.gz /app/
ADD https://example.com/file /tmp/
\end{lstlisting}
    \vspace{0.15cm}
    \begin{itemize}
      \item Además de copiar, descomprime automáticamente archivos \texttt{.tar.gz}.
      \item Puede descargar desde URLs.
      \item Úsalo solo cuando necesites la funcionalidad extra.
    \end{itemize}
  \end{columns}
  \vspace{-0.05cm}
  \begin{alertblock}{Regla práctica}
    Prefiere \texttt{COPY}. Usa \texttt{ADD} solo si necesitas descomprimir automáticamente.
  \end{alertblock}
\end{frame}

\begin{frame}[fragile]{WORKDIR: el directorio de trabajo}
  \footnotesize
  \begin{columns}[T]
    \column{0.52\textwidth}
\begin{lstlisting}[language=Dockerfile]
WORKDIR /app
\end{lstlisting}
    \vspace{0.2cm}
    \begin{itemize}
      \item Establece el directorio de trabajo para las instrucciones siguientes.
      \item Si la carpeta no existe, Docker la \textbf{crea automáticamente}.
      \item Afecta a \texttt{RUN}, \texttt{CMD}, \texttt{ENTRYPOINT}, \texttt{COPY} y \texttt{ADD}.
    \end{itemize}
    \column{0.44\textwidth}
    \begin{block}{¿Por qué usarlo?}
      \begin{itemize}
        \item Evita rutas absolutas en cada comando.
        \item Hace el Dockerfile más legible.
        \item Previene errores de archivos en carpetas equivocadas.
      \end{itemize}
    \end{block}
    \begin{alertblock}{Prefiere WORKDIR}
      Prefiere \texttt{WORKDIR /app}; \texttt{RUN cd ...} no persiste entre instrucciones.
    \end{alertblock}
  \end{columns}
\end{frame}

\begin{frame}[fragile]{ENV: variables de entorno}
  \footnotesize
  \begin{columns}[T]
    \column{0.54\textwidth}
\begin{lstlisting}[language=Dockerfile]
ENV FLASK_ENV=production
ENV APP_HOME=/app
ENV PYTHONUNBUFFERED=1
\end{lstlisting}
    \vspace{0.2cm}
    \begin{itemize}
      \item Define variables de entorno disponibles en tiempo de build y en el contenedor.
      \item Ideal para configuraciones que no cambian entre entornos.
      \item Se pueden sobrescribir con \texttt{docker run -e}.
    \end{itemize}
    \column{0.42\textwidth}
    \begin{exampleblock}{Cuándo usar ENV}
      \begin{itemize}
        \item Configuración por defecto de la app.
        \item Versiones de herramientas.
        \item Rutas y paths estándar.
      \end{itemize}
    \end{exampleblock}
    \begin{alertblock}{Cuándo NO}
      Secretos como contraseñas o tokens no deben quedar hardcodeados en la imagen.
    \end{alertblock}
  \end{columns}
\end{frame}

\begin{frame}[fragile]{EXPOSE: declarar puertos}
  \begin{columns}[T]
    \column{0.50\textwidth}
\begin{lstlisting}[language=Dockerfile]
EXPOSE 5000
EXPOSE 80/tcp
\end{lstlisting}
    \vspace{0.2cm}
    \begin{itemize}
      \item \textbf{Documenta} qué puertos escucha el contenedor.
      \item No publica el puerto automáticamente.
      \item Sirve como metadata para otros desarrolladores y herramientas.
    \end{itemize}
    \column{0.46\textwidth}
    \begin{block}{EXPOSE vs -p}
      \begin{itemize}
        \item \texttt{EXPOSE}: declara (informativo).
        \item \texttt{docker run -p}: publica y mapea el puerto.
      \end{itemize}
    \end{block}
    \vspace{0.2cm}
    \begin{center}
    \begin{tikzpicture}[
      box/.style={draw, rounded corners=4pt, minimum width=2.2cm, minimum height=0.9cm, align=center, font=\scriptsize},
      arr/.style={-{Stealth[length=2mm]}, thick}
    ]
      \node[box, fill=blue!10, draw=blue!60] (exp) {EXPOSE 5000\\\small Solo declara};
      \node[box, fill=green!12, draw=green!60!black, right=0.5cm of exp] (p) {-p 5000:5000\\\small Publica el puerto};
      \draw[arr, OTIred] (exp) -- node[above, font=\tiny] {+} (p);
    \end{tikzpicture}
    \end{center}
  \end{columns}
\end{frame}

\begin{frame}[fragile]{CMD vs ENTRYPOINT}
  \footnotesize
  \begin{columns}[T]
    \column{0.50\textwidth}
    \begin{block}{CMD}
\begin{lstlisting}[language=Dockerfile]
CMD ["python", "app.py"]
\end{lstlisting}
      \begin{itemize}
        \item Comando por defecto al iniciar el contenedor.
        \item Se puede \textbf{sobrescribir} desde \texttt{docker run}.
        \item Forma exec (array) recomendada sobre forma shell.
      \end{itemize}
    \end{block}
    \column{0.46\textwidth}
    \begin{block}{ENTRYPOINT}
\begin{lstlisting}[language=Dockerfile]
ENTRYPOINT ["python"]
CMD ["app.py"]
\end{lstlisting}
      \begin{itemize}
        \item Define el ejecutable principal (no se sobrescribe fácilmente).
        \item CMD proporciona argumentos por defecto al ENTRYPOINT.
        \item Útil para crear herramientas CLI en contenedores.
      \end{itemize}
    \end{block}
  \end{columns}
  \vspace{-0.05cm}
  \begin{exampleblock}{Regla general}
    Usa \texttt{CMD} para apps. Usa \texttt{ENTRYPOINT + CMD} para herramientas CLI.
  \end{exampleblock}
\end{frame}

\begin{frame}{Resumen visual de instrucciones}
  \scriptsize
  \begin{tabular}{p{0.15\textwidth} p{0.35\textwidth} p{0.42\textwidth}}
    \toprule
    Instrucción & ¿Qué hace? & ¿Cuándo usarla? \\
    \midrule
    \texttt{FROM} & Define la imagen base & Siempre, como primera línea \\
    \texttt{RUN} & Ejecuta comandos en el build & Instalar dependencias, scripts de setup \\
    \texttt{COPY} & Copia archivos locales a la imagen & Mover código y archivos del proyecto \\
    \texttt{WORKDIR} & Establece el directorio de trabajo & Al inicio y cuando cambies de contexto \\
    \texttt{ENV} & Define variables de entorno & Configuraciones por defecto \\
    \texttt{EXPOSE} & Declara puertos de escucha & Documentar qué puertos usa la app \\
    \texttt{CMD} & Comando por defecto del contenedor & Iniciar la aplicación \\
    \bottomrule
  \end{tabular}
\end{frame}

\MainHeaderSlide{Capas, caché y\\buenas prácticas}{Cómo Docker construye imágenes y cómo optimizarlas}
\section{S2 · Capas, Caché y Buenas Prácticas}

\begin{frame}{¿Cómo construye Docker una imagen?}
  \begin{center}
    \begin{tikzpicture}[
      layer/.style={draw, rounded corners=4pt, minimum width=7cm, minimum height=0.70cm, align=center, font=\scriptsize},
    ]
      \node[layer, fill=blue!10, draw=blue!60] (l0) at (0,0) {FROM python:3.12-slim \hspace{2cm} \textit{← capa de la imagen base}};
      \node[layer, fill=green!8, draw=green!50!black, below=0.04cm of l0] (l1) {WORKDIR /app \hspace{2.65cm} \textit{← capa 1}};
      \node[layer, fill=orange!10, draw=orange!70, below=0.04cm of l1] (l2) {COPY requirements.txt . \hspace{2.1cm} \textit{← capa 2}};
      \node[layer, fill=green!8, draw=green!50!black, below=0.04cm of l2] (l3) {RUN pip install -r requirements.txt \hspace{1.2cm} \textit{← capa 3}};
      \node[layer, fill=orange!10, draw=orange!70, below=0.04cm of l3] (l4) {COPY app.py . \hspace{3.3cm} \textit{← capa 4}};
      \node[layer, fill=OTIred!10, draw=OTIred!70, below=0.04cm of l4] (l5) {CMD ["python", "app.py"] \hspace{2.2cm} \textit{← capa 5 (metadatos)}};
    \end{tikzpicture}
  \end{center}
  \vspace{-0.05cm}
  \begin{block}{Cada instrucción del Dockerfile genera una capa}
    Las capas son \textbf{inmutables}. Docker reutiliza las que no cambiaron mediante caché.
  \end{block}
\end{frame}

\begin{frame}{La caché de Docker: cómo funciona}
  \begin{columns}[c]
    \column{0.56\textwidth}
    \footnotesize
    \begin{itemize}
      \item Docker compara cada instrucción con la caché de builds anteriores.
      \item Si la instrucción y sus dependencias no cambiaron, \textbf{reutiliza la capa}.
      \item Si una capa cambia, \textbf{todas las capas siguientes se reconstruyen}.
    \end{itemize}
    \vspace{0.15cm}
    \begin{alertblock}{Regla de oro del orden}
      Pon primero lo que \textbf{cambia menos} (instalación de dependencias) y al final lo que \textbf{cambia más} (código de la app).
    \end{alertblock}
    \column{0.40\textwidth}
    \centering
    \begin{tikzpicture}[
      box/.style={draw, rounded corners=4pt, minimum width=3cm, minimum height=0.8cm, align=center, font=\scriptsize},
    ]
      \node[box, fill=green!12, draw=green!60!black] (a) {COPY requirements.txt\\\small \textit{(cambia poco)}};
      \node[box, fill=green!12, draw=green!60!black, below=0.3cm of a] (b) {RUN pip install\\\small \textit{(usa caché)}};
      \node[box, fill=orange!15, draw=orange!70, below=0.3cm of b] (c) {COPY app.py\\\small \textit{(cambia mucho)}};
      \draw[{Stealth[length=3mm]}-{Stealth[length=3mm]}, thick, OTIred] (a.east) -- ++(0.8,0) node[right, font=\tiny] {Cache};
      \draw[{Stealth[length=3mm]}-{Stealth[length=3mm]}, thick, OTIred] (b.east) -- ++(0.8,0) node[right, font=\tiny] {Cache};
      \node[font=\tiny, color=UUdarkgray, right=0.3cm of c.east] {Reconstruye};
    \end{tikzpicture}
  \end{columns}
\end{frame}

\begin{frame}[fragile]{Ejemplo: Dockerfile optimizado}
  \begin{columns}[T]
    \column{0.48\textwidth}
    \begin{alertblock}{Mal orden (lento)}
\begin{lstlisting}[language=Dockerfile, basicstyle=\ttfamily\tiny]
FROM python:3.12

COPY . /app          # TODO cambia
WORKDIR /app

RUN pip install \    # Siempre
    -r requirements.txt

CMD ["python","app.py"]
\end{lstlisting}
    \end{alertblock}
    \column{0.48\textwidth}
    \begin{exampleblock}{Buen orden (rápido)}
\begin{lstlisting}[language=Dockerfile, basicstyle=\ttfamily\tiny]
FROM python:3.12-slim

WORKDIR /app

# Dependencias primero
COPY requirements.txt .
RUN pip install --no-cache-dir \
    -r requirements.txt

# Codigo al final
COPY . .

CMD ["python","app.py"]
\end{lstlisting}
    \end{exampleblock}
  \end{columns}
  \vspace{-0.05cm}
  \begin{block}{¿Por qué es más rápido?}
    Cachea dependencias; cambios en \texttt{app.py} no reinstalan paquetes.
  \end{block}
\end{frame}

\begin{frame}[fragile]{Otras buenas prácticas}
  \footnotesize
  \begin{columns}[T]
    \column{0.48\textwidth}
    \begin{block}{1. Usa imágenes base específicas}
\begin{lstlisting}[language=Dockerfile, basicstyle=\ttfamily\tiny]
# Mal: imagen generica y pesada
FROM python

# Bien: version fija y liviana
FROM python:3.12-slim
\end{lstlisting}
    \end{block}
    \begin{block}{2. No instales paquetes innecesarios}
\begin{lstlisting}[language=Dockerfile, basicstyle=\ttfamily\tiny]
RUN apt-get update && \
    apt-get install -y --no-install-recommends \
    curl && \
    rm -rf /var/lib/apt/lists/*
\end{lstlisting}
    \end{block}
    \column{0.48\textwidth}
    \begin{block}{3. Limpia en el mismo RUN}
\begin{lstlisting}[language=Dockerfile, basicstyle=\ttfamily\tiny]
RUN pip install -r requirements.txt && \
    rm -rf /root/.cache/pip
\end{lstlisting}
    \end{block}
    \begin{block}{4. Usa multistage builds}
\begin{lstlisting}[language=Dockerfile, basicstyle=\ttfamily\tiny]
# Stage 1: build
FROM golang:1.21 AS builder
RUN go build -o app

# Stage 2: runtime (imagen final pequeña)
FROM alpine:3.19
COPY --from=builder /app .
CMD ["./app"]
\end{lstlisting}
    \end{block}
  \end{columns}
\end{frame}

\begin{frame}{Multistage build: el patrón profesional}
  \begin{center}
    \begin{tikzpicture}[
      stage/.style={draw, rounded corners=6pt, minimum width=4cm, minimum height=2cm, align=center, font=\small},
      arr/.style={-{Stealth[length=4mm]}, very thick, dockerblue}
    ]
      \node[stage, fill=orange!15, draw=orange!70] (build) at (0,0) {
        {\faCog\ \textbf{Stage 1: Builder}}\\[4pt]
        \scriptsize
        Imagen completa con\\compiladores y herramientas
      };
      \node[stage, fill=green!12, draw=green!60!black, right=1.5cm of build] (prod) {
        {\faRocket\ \textbf{Stage 2: Runtime}}\\[4pt]
        \scriptsize
        Solo el binario,\\imagen mínima
      };
      \draw[arr] (build) -- node[above, font=\scriptsize] {COPY --from=builder} (prod);

      \node[font=\footnotesize, below=0.8cm of build, text width=4cm, align=center] {\color{UUdarkgray}~500 MB};
      \node[font=\footnotesize, below=0.8cm of prod, text width=4cm, align=center] {\color{UUdarkgray}~15 MB};
    \end{tikzpicture}
  \end{center}
  \vspace{0.2cm}
  \begin{block}{Ventaja}
    La imagen final solo contiene lo necesario para ejecutar. Las herramientas de compilación quedan en la etapa de build y no contaminan la imagen de producción.
  \end{block}
\end{frame}

\MainHeaderSlide{.dockerignore e\\imágenes base livianas}{Evitar basura en la imagen y elegir la base correcta}
\section{S2 · .dockerignore e Imágenes Base Livianas}

\begin{frame}[fragile]{.dockerignore: qué es y para qué sirve}
  \footnotesize
  \begin{columns}[c]
    \column{0.52\textwidth}
\begin{lstlisting}[style=terminal, basicstyle=\ttfamily\scriptsize]
# .dockerignore
__pycache__
*.pyc
*.pyo
.git
.env
*.log
node_modules
venv/
.venv/
.vscode/
*.md
Dockerfile
docker-compose.yml
\end{lstlisting}
    \column{0.44\textwidth}
    \begin{itemize}
      \item Similar a \texttt{.gitignore}: archivos y carpetas que \textbf{no} se envían al contexto de build.
      \item Reduce el tamaño del contexto y acelera el build.
      \item Evita que archivos sensibles (\texttt{.env}, claves) terminen en la imagen.
    \end{itemize}
  \end{columns}
  \vspace{-0.05cm}
  \begin{alertblock}{Regla}
    Incluye \texttt{.dockerignore}; evita enviar \texttt{node\_modules}, \texttt{.git} y credenciales.
  \end{alertblock}
\end{frame}

\begin{frame}{Imágenes base: slim vs alpine}
  \begin{center}
    \begin{tikzpicture}[
      box/.style={draw, rounded corners=6pt, minimum width=3.8cm, minimum height=2.2cm, align=center, font=\small},
    ]
      \node[box, fill=OTIred!10, draw=OTIred!70] (full) at (0,0) {
        {\color{OTIred}\faWeightHanging\ \textbf{Completa}}\\[4pt]
        \texttt{python:3.12}\\[3pt]
        \scriptsize
        ~1 GB\\Todas las herramientas\\Debian completo
      };
      \node[box, fill=orange!15, draw=orange!70, right=0.8cm of full] (slim) {
        {\color{orange!80}\faBalanceScale\ \textbf{Slim}}\\[4pt]
        \texttt{python:3.12-slim}\\[3pt]
        \scriptsize
        ~150 MB\\Debian mínimo\\Lo esencial para Python
      };
      \node[box, fill=green!12, draw=green!60!black, right=0.8cm of slim] (alpine) {
        {\color{green!50!black}\faFeather\ \textbf{Alpine}}\\[4pt]
        \texttt{python:3.12-alpine}\\[3pt]
        \scriptsize
        ~50 MB\\Alpine Linux (musl)\\La más liviana
      };
    \end{tikzpicture}
  \end{center}
  \vspace{0.2cm}
  \begin{block}{¿Cuál elegir?}
    \textbf{Slim}: buena relación tamaño/compatibilidad. \textbf{Alpine}: máxima reducción, pero puede tener problemas con paquetes que dependen de glibc. Para empezar, \texttt{slim} es la opción más segura.
  \end{block}
\end{frame}

\MainHeaderSlide{Publicación en\\Docker Hub}{Tagging, versionado y subida de imágenes}
\section{S2 · Publicación en Docker Hub}

\begin{frame}{Flujo de publicación}
  \begin{center}
    \begin{tikzpicture}[
      stage/.style={draw, rounded corners=6pt, minimum width=2.5cm, minimum height=1.3cm, align=center, font=\small},
      arr/.style={-{Stealth[length=3mm]}, thick, dockerblue}
    ]
      \node[stage, fill=orange!15, draw=orange!70] (login) at (0,0) {\faUser\\\textbf{Login}\\\texttt{docker login}};
      \node[stage, fill=blue!10, draw=blue!60, right=0.7cm of login] (tag) {\faTag\\\textbf{Etiquetar}\\\texttt{docker tag}};
      \node[stage, fill=green!12, draw=green!60!black, right=0.7cm of tag] (push) {\faUpload\\\textbf{Publicar}\\\texttt{docker push}};
      \node[stage, fill=OTIred!10, draw=OTIred!70, right=0.7cm of push] (pull) {\faDownload\\\textbf{Descargar}\\\texttt{docker pull}};
      \draw[arr] (login) -- (tag);
      \draw[arr] (tag) -- (push);
      \draw[arr] (push) -- (pull);
    \end{tikzpicture}
  \end{center}
  \vspace{0.2cm}
  \begin{block}{El ciclo completo}
    1. Construyes tu imagen. \quad 2. Le asignas un tag. \quad 3. La publicas en un registro. \quad 4. Cualquiera puede descargarla y usarla.
  \end{block}
\end{frame}

\begin{frame}[fragile]{Login y tagging}
  \scriptsize
  \begin{columns}[T]
    \column{0.52\textwidth}
    \begin{block}{1. Iniciar sesión en Docker Hub}
\begin{lstlisting}[style=terminal, basicstyle=\ttfamily\tiny]
docker login
# Usuario: tu-usuario
# Password: tu-token
\end{lstlisting}
    \end{block}
    \vspace{-0.05cm}
    \begin{block}{2. Etiquetar la imagen}
\begin{lstlisting}[style=terminal, basicstyle=\ttfamily\tiny]
# Formato:
# usuario/imagen:tag

docker tag mi-flask:v1 \
  tony/mi-flask:v1
\end{lstlisting}
    \end{block}
    \column{0.44\textwidth}
    \begin{exampleblock}{Convención de tags}
      \begin{itemize}
        \item \texttt{v1, v2, v3}: versiones específicas.
        \item \texttt{latest}: la más reciente (se actualiza manualmente).
        \item \texttt{dev, staging, prod}: por entorno.
        \item \texttt{2026-05-31}: por fecha de build.
      \end{itemize}
    \end{exampleblock}
    \begin{alertblock}{Importante}
      \texttt{latest} no es automático: tú decides qué imagen lo lleva.
    \end{alertblock}
  \end{columns}
\end{frame}

\begin{frame}[fragile]{Push y pull}
  \footnotesize
  \begin{columns}[T]
    \column{0.50\textwidth}
    \begin{block}{3. Publicar la imagen}
\begin{lstlisting}[style=terminal, basicstyle=\ttfamily\scriptsize]
docker push tony/mi-flask:v1
docker push tony/mi-flask:latest
\end{lstlisting}
    \vspace{0.1cm}
    La imagen queda disponible en Docker Hub según la visibilidad del repositorio.
    \end{block}
    \column{0.46\textwidth}
    \begin{block}{4. Descargar y ejecutar}
\begin{lstlisting}[style=terminal, basicstyle=\ttfamily\scriptsize]
docker pull tony/mi-flask:v1

docker run -d -p 5000:5000 \
  tony/mi-flask:v1
\end{lstlisting}
    \end{block}
    \vspace{0.15cm}
    \begin{exampleblock}{¿Pública o privada?}
      Docker Hub permite 1 repo privado gratis. GHCR es alternativa para repos privados.
    \end{exampleblock}
  \end{columns}
\end{frame}

\MainHeaderSlide{Laboratorio\\práctico}{Construir, optimizar y publicar una imagen profesional}
\section{S2 · Laboratorio Práctico}

\begin{frame}{Laboratorio: app Flask optimizada}
  \begin{center}
    \begin{tikzpicture}[
      stepnode/.style={draw, rounded corners=6pt, minimum width=2.5cm, minimum height=1.2cm, align=center, font=\scriptsize},
      arr/.style={-{Stealth[length=3mm]}, thick, OTIred}
    ]
      \node[stepnode, fill=orange!15, draw=orange!70] (s1) {\faFileCode\\\textbf{Escribir}\\Dockerfile optimizado};
      \node[stepnode, fill=blue!10, draw=blue!60, right=0.5cm of s1] (s2) {\faCog\\\textbf{Build}\\Aprovechar caché};
      \node[stepnode, fill=green!12, draw=green!60!black, right=0.5cm of s2] (s3) {\faPlay\\\textbf{Ejecutar}\\Probar la app};
      \node[stepnode, fill=OTIred!10, draw=OTIred!70, right=0.5cm of s3] (s4) {\faUpload\\\textbf{Publicar}\\Docker Hub};
      \draw[arr] (s1) -- (s2);
      \draw[arr] (s2) -- (s3);
      \draw[arr] (s3) -- (s4);
    \end{tikzpicture}
  \end{center}
  \vspace{0.2cm}
  \begin{block}{Objetivo}
    Partiendo del proyecto Flask de la sesión 1, optimizar el Dockerfile con buenas prácticas, construir la imagen aprovechando caché y publicarla en Docker Hub.
  \end{block}
\end{frame}

\begin{frame}[fragile]{Dockerfile optimizado para el laboratorio}
\begin{lstlisting}[language=Dockerfile, basicstyle=\ttfamily\scriptsize]
# Etapa 1: builder (si necesitas compilar)
FROM python:3.12-slim AS builder
WORKDIR /app
COPY requirements.txt .
RUN pip install --no-cache-dir -r requirements.txt

# Etapa 2: imagen final (runtime)
FROM python:3.12-slim
WORKDIR /app
# Copiar dependencias ya instaladas desde el builder
COPY --from=builder /usr/local/lib/python3.12/site-packages \
     /usr/local/lib/python3.12/site-packages
COPY --from=builder /usr/local/bin /usr/local/bin
COPY app.py .
EXPOSE 5000
CMD ["python", "app.py"]
\end{lstlisting}
\end{frame}

\begin{frame}[fragile]{Comandos del laboratorio}
  \footnotesize
  \begin{columns}[T]
    \column{0.50\textwidth}
\begin{lstlisting}[style=terminal]
# 1. Construir la imagen
docker build -t mi-flask:v2 .

# 2. Ver las capas generadas
docker history mi-flask:v2

# 3. Ejecutar y probar
docker run -d --name flask-v2 \
  -p 5000:5000 mi-flask:v2

curl http://localhost:5000
\end{lstlisting}
    \column{0.46\textwidth}
\begin{lstlisting}[style=terminal]
# 4. Login en Docker Hub
docker login

# 5. Etiquetar
docker tag mi-flask:v2 \
  tony/mi-flask:v2

# 6. Publicar
docker push tony/mi-flask:v2

# 7. Verificar en:
# hub.docker.com/r/tony/mi-flask
\end{lstlisting}
    \end{columns}
\end{frame}

\begin{frame}{Errores frecuentes en esta sesión}
  \scriptsize
  \begin{tabular}{p{0.28\textwidth} p{0.32\textwidth} p{0.32\textwidth}}
    \toprule
    Error & Causa probable & Solución rápida \\
    \midrule
    Build lento siempre & Mal orden de capas (COPY . . antes de RUN) & Mover dependencias antes del código \\
    Imagen muy pesada & Usar imagen base \texttt{:latest} completa & Cambiar a \texttt{:slim} o \texttt{:alpine} \\
    Caché no funciona & Cambió un archivo antes de lo esperado & Revisar orden: lo estable primero \\
    \texttt{docker push} rechazado & No hay login o tag mal escrito & \texttt{docker login} y verificar \texttt{usuario/imagen:tag} \\
    Archivos \texttt{.env} en la imagen & Falta \texttt{.dockerignore} & Crear \texttt{.dockerignore} y reconstruir \\
    Error \texttt{no space left} & Imágenes acumuladas en disco & \texttt{docker system prune -a} \\
    \bottomrule
  \end{tabular}
\end{frame}

\begin{frame}{Checklist de aprendizaje — Sesión 2}
  \begin{itemize}
    \item Puedo escribir un Dockerfile con FROM, RUN, COPY, WORKDIR, ENV, EXPOSE y CMD.
    \item Puedo explicar la diferencia entre COPY y ADD, CMD y ENTRYPOINT.
    \item Puedo explicar qué es una capa y cómo afecta el orden de las instrucciones.
    \item Puedo optimizar un Dockerfile para aprovechar la caché.
    \item Puedo usar .dockerignore para excluir archivos del contexto de build.
    \item Puedo elegir entre imagen base completa, slim y alpine.
    \item Puedo publicar una imagen en Docker Hub con tagging y versionado.
    \item Puedo descargar y ejecutar una imagen desde Docker Hub.
  \end{itemize}
\end{frame}

\begin{frame}{Resumen de la Sesión 2}
  \footnotesize
  \begin{enumerate}
    \item Un Dockerfile define la receta para construir imágenes propias.
    \item FROM, RUN, COPY, ENV, EXPOSE y CMD cubren el flujo base de una app.
    \item El orden de capas afecta directamente la velocidad del build.
    \item .dockerignore evita enviar basura o secretos al contexto de build.
    \item Slim y Alpine ayudan a reducir tamaño, con criterios distintos.
    \item Docker Hub permite compartir imágenes mediante tags versionados.
  \end{enumerate}
  \vspace{0.08cm}
  \begin{alertblock}{Ahora: Sesión 3}
    \textbf{Docker Compose:} levantar aplicaciones multi-contenedor con Flask, PostgreSQL y variables de entorno.
  \end{alertblock}
\end{frame}

% ============================================================
% SESIÓN 3: DOCKER COMPOSE Y APLICACIONES MULTI-CONTENEDOR
% ============================================================

\begin{frame}{Objetivo de la sesión 3}
  \footnotesize
  \begin{columns}[c]
    \column{0.58\textwidth}
    \begin{block}{Al terminar la sesión 3 podrás}
      \begin{itemize}
        \item Explicar para qué sirve Docker Compose.
        \item Levantar proyectos reales con múltiples servicios.
        \item Entender la estructura de \texttt{docker-compose.yml}.
        \item Conectar una aplicación Flask con PostgreSQL.
        \item Configurar variables de entorno usando archivo \texttt{.env}.
        \item Ejecutar, revisar y detener un stack completo con Compose.
      \end{itemize}
    \end{block}
    \column{0.36\textwidth}
    \centering
    {\color{dockerblue}\fontsize{42}{48}\selectfont\faLayerGroup}\\[0.15cm]
    {\small\textbf{Un comando, varios servicios}}
  \end{columns}
\end{frame}

\MainHeaderSlide{Docker Compose}{Levantar proyectos reales con múltiples servicios}
\section{S3 · Docker Compose}

\begin{frame}{El problema: una app ya no vive sola}
  \begin{columns}[c]
    \column{0.55\textwidth}
    \small
    \begin{itemize}
      \item Una app real suele necesitar base de datos, cache, proxy o workers.
      \item Ejecutar cada contenedor a mano vuelve el laboratorio repetitivo.
      \item Los comandos \texttt{docker run} crecen con puertos, redes, volúmenes y variables.
      \item El equipo necesita una forma reproducible de levantar todo el stack.
    \end{itemize}
    \begin{alertblock}{Pregunta clave}
      ¿Cómo describimos una aplicación completa, no solo un contenedor?
    \end{alertblock}
    \column{0.39\textwidth}
    \centering
    {\color{OTIred}\fontsize{46}{52}\selectfont\faProjectDiagram}\\[0.15cm]
    {\large\textbf{Multi-servicio}}
  \end{columns}
\end{frame}

\begin{frame}{Qué es Docker Compose}
  \begin{center}
  \begin{tikzpicture}[
    box/.style={draw, rounded corners=6pt, minimum width=3.2cm, minimum height=1.25cm, align=center, text width=3.0cm, font=\small},
    arr/.style={-{Stealth[length=3mm]}, thick, dockerblue}
  ]
    \node[box, fill=orange!15, draw=orange!70] (yaml) {\faFileCode\\\textbf{compose.yml}\\Servicios declarados};
    \node[box, fill=blue!10, draw=blue!60, right=0.8cm of yaml] (compose) {\faLayerGroup\\\textbf{Docker Compose}\\Orquesta localmente};
    \node[box, fill=green!12, draw=green!60!black, right=0.8cm of compose] (stack) {\faCubes\\\textbf{Stack}\\App + BD + red};
    \draw[arr] (yaml) -- (compose);
    \draw[arr] (compose) -- (stack);
  \end{tikzpicture}
  \end{center}
  \vspace{0.2cm}
  \begin{block}{Definición práctica}
    Docker Compose permite definir y ejecutar aplicaciones multi-contenedor usando un archivo YAML versionable.
  \end{block}
\end{frame}

\begin{frame}[fragile]{De docker run a docker compose}
  \footnotesize
  \begin{columns}[T]
    \column{0.48\textwidth}
    \begin{alertblock}{Sin Compose}
\begin{lstlisting}[style=terminal, basicstyle=\ttfamily\tiny]
docker network create app-net
docker run -d --name db \
  --network app-net postgres:16
docker run -d --name web \
  --network app-net -p 5000:5000 \
  mi-flask:v3
\end{lstlisting}
    \end{alertblock}
    \column{0.48\textwidth}
    \begin{exampleblock}{Con Compose}
\begin{lstlisting}[style=terminal, basicstyle=\ttfamily\tiny]
docker compose up -d
docker compose ps
docker compose down
\end{lstlisting}
      \vspace{0.1cm}
      El archivo \texttt{docker-compose.yml} guarda la arquitectura del proyecto.
    \end{exampleblock}
  \end{columns}
\end{frame}

\MainHeaderSlide{Estructura de compose.yml}{Servicios, redes, volúmenes y dependencias}
\section{S3 · Estructura de docker-compose.yml}

\begin{frame}{Piezas principales de Compose}
  \begin{center}
  \begin{tikzpicture}[
    box/.style={draw, rounded corners=6pt, minimum width=4.6cm, minimum height=1.45cm, align=center, text width=4.3cm, font=\small},
  ]
    \node[box, fill=blue!10, draw=blue!60] (services) at (0,0) {\faCubes\ \textbf{services}\\Contenedores del proyecto};
    \node[box, fill=green!12, draw=green!60!black, right=0.7cm of services] (networks) {\faNetworkWired\ \textbf{networks}\\Comunicación interna};
    \node[box, fill=orange!15, draw=orange!70, below=0.45cm of services] (volumes) {\faDatabase\ \textbf{volumes}\\Datos persistentes};
    \node[box, fill=OTIred!10, draw=OTIred!70, below=0.45cm of networks] (depends) {\faLink\ \textbf{depends\_on}\\Orden de arranque};
  \end{tikzpicture}
  \end{center}
  \vspace{0.1cm}
  \begin{block}{Idea mental}
    Compose describe la topología local: qué servicios existen, cómo se conectan y qué datos conservan.
  \end{block}
\end{frame}

\begin{frame}[fragile]{docker-compose.yml mínimo}
\begin{lstlisting}[style=terminal, basicstyle=\ttfamily\tiny]
services:
  web:
    build: .
    ports:
      - "5000:5000"
    depends_on:
      - db
  db:
    image: postgres:16
    environment:
      POSTGRES_DB: appdb
      POSTGRES_USER: appuser
      POSTGRES_PASSWORD: apppass
\end{lstlisting}
  \vspace{-0.08cm}
  \begin{alertblock}{Lectura rápida}
    \texttt{web} construye la app. \texttt{db} usa PostgreSQL. Compose crea una red interna por defecto.
  \end{alertblock}
\end{frame}

\begin{frame}[fragile]{Comandos esenciales de Compose}
  \scriptsize
  \begin{tabular}{p{0.30\textwidth} p{0.34\textwidth} p{0.25\textwidth}}
    \toprule
    Comando & Para qué sirve & Uso típico \\
    \midrule
    \texttt{docker compose up -d} & Levanta el stack en segundo plano & Iniciar laboratorio \\
    \texttt{docker compose ps} & Lista servicios del proyecto & Ver estado \\
    \texttt{docker compose logs -f} & Muestra logs integrados & Diagnóstico \\
    \texttt{docker compose exec web sh} & Abre shell en un servicio & Debug \\
    \texttt{docker compose down} & Detiene y elimina contenedores/red & Limpiar stack \\
    \bottomrule
  \end{tabular}
\end{frame}

\begin{frame}{Dependencias entre servicios}
  \begin{columns}[c]
    \column{0.52\textwidth}
    \small
    \begin{itemize}
      \item \texttt{depends\_on} define orden de arranque básico.
      \item No garantiza que la base de datos ya acepte conexiones.
      \item Para producción o laboratorios robustos se usan healthchecks.
      \item En esta sesión nos enfocamos en levantar y conectar el stack.
    \end{itemize}
    \begin{alertblock}{Regla práctica}
      Arrancar antes no significa estar listo para recibir tráfico.
    \end{alertblock}
    \column{0.42\textwidth}
    \centering
    \begin{tikzpicture}[
      box/.style={draw, rounded corners=6pt, minimum width=3.2cm, minimum height=1.1cm, align=center, font=\small},
      arr/.style={-{Stealth[length=3mm]}, thick, OTIred}
    ]
      \node[box, fill=green!12, draw=green!60!black] (db) {\faDatabase\\\textbf{db}};
      \node[box, fill=blue!10, draw=blue!60, above=0.65cm of db] (web) {\faPython\\\textbf{web}};
      \draw[arr] (web) -- node[right, font=\scriptsize] {depends\_on} (db);
    \end{tikzpicture}
  \end{columns}
\end{frame}

\MainHeaderSlide{Flask + PostgreSQL}{Integrar una aplicación con base de datos}
\section{S3 · Flask con PostgreSQL}

\begin{frame}{Arquitectura del laboratorio}
  \begin{center}
  \begin{tikzpicture}[
    box/.style={draw, rounded corners=6pt, minimum width=3.2cm, minimum height=1.3cm, align=center, text width=3.0cm, font=\small},
    arr/.style={-{Stealth[length=3mm]}, thick, dockerblue}
  ]
    \node[box, fill=UUlightgray, draw=UUdarkgray] (browser) {\faDesktop\\\textbf{Navegador}\\localhost:5000};
    \node[box, fill=blue!10, draw=blue!60, right=0.8cm of browser] (web) {\faPython\\\textbf{Flask}\\servicio \texttt{web}};
    \node[box, fill=green!12, draw=green!60!black, right=0.8cm of web] (db) {\faDatabase\\\textbf{PostgreSQL}\\servicio \texttt{db}};
    \draw[arr] (browser) -- node[above, font=\scriptsize] {HTTP} (web);
    \draw[arr] (web) -- node[above, font=\scriptsize] {db:5432} (db);
  \end{tikzpicture}
  \end{center}
  \vspace{0.15cm}
  \begin{block}{Clave de red interna}
    Dentro de Compose, la app no se conecta a \texttt{localhost}; se conecta al nombre del servicio: \texttt{db}.
  \end{block}
\end{frame}

\begin{frame}[fragile]{Variables de conexión en Flask}
  \footnotesize
  \begin{columns}[T]
    \column{0.50\textwidth}
\begin{lstlisting}[language=Python, basicstyle=\ttfamily\scriptsize]
import os

DB_HOST = os.getenv("DB_HOST", "db")
DB_NAME = os.getenv("DB_NAME", "appdb")
DB_USER = os.getenv("DB_USER", "appuser")
DB_PASSWORD = os.getenv("DB_PASSWORD")
\end{lstlisting}
    \column{0.46\textwidth}
    \begin{block}{Por qué usar variables}
      \begin{itemize}
        \item Evitan hardcodear credenciales.
        \item Permiten cambiar entorno sin tocar código.
        \item Funcionan bien con \texttt{.env} y Compose.
      \end{itemize}
    \end{block}
  \end{columns}
\end{frame}

\begin{frame}[fragile]{Compose para Flask + PostgreSQL}
\begin{lstlisting}[style=terminal, basicstyle=\ttfamily\tiny]
services:
  web:
    build: .
    ports:
      - "5000:5000"
    env_file:
      - .env
    depends_on:
      - db

  db:
    image: postgres:16
    environment:
      POSTGRES_DB: ${DB_NAME}
      POSTGRES_USER: ${DB_USER}
      POSTGRES_PASSWORD: ${DB_PASSWORD}
    volumes:
      - postgres_data:/var/lib/postgresql/data

volumes:
  postgres_data:
\end{lstlisting}
\end{frame}

\MainHeaderSlide{Variables de entorno}{Configurar el proyecto con archivo .env}
\section{S3 · Variables con .env}

\begin{frame}[fragile]{Archivo .env del proyecto}
  \begin{columns}[T]
    \column{0.50\textwidth}
\begin{lstlisting}[style=terminal, basicstyle=\ttfamily\scriptsize]
DB_HOST=db
DB_NAME=appdb
DB_USER=appuser
DB_PASSWORD=appsecret
FLASK_ENV=development
\end{lstlisting}
    \column{0.45\textwidth}
    \begin{alertblock}{Importante}
      No subas \texttt{.env} con credenciales reales a GitHub. Usa \texttt{.env.example} para documentar variables esperadas.
    \end{alertblock}
    \begin{block}{Para clase}
      Usaremos valores simples de laboratorio para entender el flujo.
    \end{block}
  \end{columns}
\end{frame}

\begin{frame}{.env vs environment}
  \footnotesize
  \begin{tabular}{p{0.28\textwidth} p{0.31\textwidth} p{0.31\textwidth}}
    \toprule
    Opción & Ventaja & Cuándo usarla \\
    \midrule
    \texttt{environment} & Visible en el YAML & Valores simples o explícitos \\
    \texttt{env\_file} & Separa configuración & Variables por entorno \\
    \texttt{.env.example} & Documenta sin secretos & Repos compartidos \\
    Secret manager & Mejor seguridad & Producción real \\
    \bottomrule
  \end{tabular}
  \vspace{0.2cm}
  \begin{alertblock}{Regla práctica}
    El archivo \texttt{.env} ayuda en desarrollo; no es una bóveda de secretos.
  \end{alertblock}
\end{frame}

\MainHeaderSlide{Laboratorio\\multi-contenedor}{Levantar Flask y PostgreSQL con Compose}
\section{S3 · Laboratorio Compose}

\begin{frame}{Flujo del laboratorio}
  \begin{center}
  \begin{tikzpicture}[
    labstep/.style={draw, rounded corners=6pt, minimum width=2.45cm, minimum height=1.15cm, align=center, text width=2.25cm, font=\scriptsize},
    arr/.style={-{Stealth[length=3mm]}, thick, OTIred}
  ]
    \node[labstep, fill=orange!15, draw=orange!70] (s1) {\faFolderOpen\\\textbf{Preparar}\\archivos};
    \node[labstep, fill=blue!10, draw=blue!60, right=0.45cm of s1] (s2) {\faFileCode\\\textbf{Compose}\\YAML};
    \node[labstep, fill=green!12, draw=green!60!black, right=0.45cm of s2] (s3) {\faPlay\\\textbf{Levantar}\\up -d};
    \node[labstep, fill=OTIred!10, draw=OTIred!70, right=0.45cm of s3] (s4) {\faSearch\\\textbf{Validar}\\logs + ps};
    \draw[arr] (s1) -- (s2);
    \draw[arr] (s2) -- (s3);
    \draw[arr] (s3) -- (s4);
  \end{tikzpicture}
  \end{center}
  \begin{block}{Meta}
    Ejecutar una aplicación Flask conectada a PostgreSQL usando un solo archivo Compose.
  \end{block}
\end{frame}

\begin{frame}[fragile]{Comandos del laboratorio Compose}
  \footnotesize
  \begin{columns}[T]
    \column{0.50\textwidth}
\begin{lstlisting}[style=terminal]
# 1. Levantar el stack
docker compose up -d --build

# 2. Revisar servicios
docker compose ps

# 3. Ver logs integrados
docker compose logs -f
\end{lstlisting}
    \column{0.46\textwidth}
\begin{lstlisting}[style=terminal]
# 4. Probar la app
curl http://localhost:5000

# 5. Entrar al servicio web
docker compose exec web sh

# 6. Detener y limpiar
docker compose down
\end{lstlisting}
  \end{columns}
\end{frame}

\begin{frame}{Errores frecuentes en Compose}
  \scriptsize
  \begin{tabular}{p{0.28\textwidth} p{0.32\textwidth} p{0.32\textwidth}}
    \toprule
    Error & Causa probable & Solución rápida \\
    \midrule
    App no conecta a BD & Usa \texttt{localhost} en vez de \texttt{db} & Cambiar host al nombre del servicio \\
    Puerto ocupado & Otro proceso usa 5000 & Cambiar a \texttt{5001:5000} \\
    Variables vacías & Falta \texttt{.env} o nombre incorrecto & Revisar \texttt{env\_file} y claves \\
    BD pierde datos & No hay volumen persistente & Agregar volumen nombrado \\
    Cambios no aparecen & Imagen antigua en caché & Usar \texttt{up -d --build} \\
    \bottomrule
  \end{tabular}
\end{frame}

\begin{frame}{Checklist de aprendizaje — Sesión 3}
  \begin{itemize}
    \item Puedo explicar para qué sirve Docker Compose.
    \item Puedo leer la estructura básica de \texttt{docker-compose.yml}.
    \item Puedo definir servicios, puertos, dependencias y volúmenes.
    \item Puedo conectar Flask con PostgreSQL usando el nombre del servicio.
    \item Puedo configurar variables con \texttt{.env} y \texttt{env\_file}.
    \item Puedo levantar y detener un stack con \texttt{docker compose up} y \texttt{down}.
  \end{itemize}
\end{frame}

\begin{frame}{Resumen de la Sesión 3}
  \footnotesize
  \begin{enumerate}
    \item Compose describe aplicaciones multi-contenedor en un archivo YAML.
    \item Los servicios se comunican por nombre dentro de la red del proyecto.
    \item \texttt{ports} publica hacia la máquina host; la red interna usa nombres de servicio.
    \item \texttt{volumes} permite conservar datos de PostgreSQL entre reinicios.
    \item \texttt{.env} separa configuración del código y del YAML.
  \end{enumerate}
  \vspace{0.08cm}
  \begin{alertblock}{Próxima sesión}
    \textbf{Redes, volúmenes y persistencia:} comunicación interna, almacenamiento y datos durables.
  \end{alertblock}
\end{frame}

% ============================================================
% SESIÓN 4: REDES, VOLÚMENES Y PERSISTENCIA
% ============================================================

\begin{frame}{Objetivo de la sesión 4}
  \footnotesize
  \begin{columns}[c]
    \column{0.58\textwidth}
    \begin{block}{Al terminar la sesión 4 podrás}
      \begin{itemize}
        \item Explicar cómo se comunican contenedores mediante redes Docker.
        \item Diferenciar red interna, puertos publicados y nombres de servicio.
        \item Usar volúmenes nombrados y bind mounts para persistencia.
        \item Aislar servicios internos como bases de datos.
        \item Agregar healthchecks para validar disponibilidad.
        \item Realizar backup y restauración básica de datos del proyecto.
      \end{itemize}
    \end{block}
    \column{0.36\textwidth}
    \centering
    {\color{dockerblue}\fontsize{42}{48}\selectfont\faNetworkWired}\\[0.15cm]
    {\small\textbf{Datos durables y redes claras}}
  \end{columns}
\end{frame}

\MainHeaderSlide{Redes Docker}{Comunicación entre contenedores}
\section{S4 · Redes Docker}

\begin{frame}{Qué problema resuelven las redes}
  \begin{columns}[c]
    \column{0.56\textwidth}
    \small
    \begin{itemize}
      \item Los contenedores necesitan comunicarse sin depender de IPs manuales.
      \item Compose crea una red interna por proyecto automáticamente.
      \item Cada servicio puede resolverse por su nombre: \texttt{web}, \texttt{db}, \texttt{cache}.
      \item Solo publicamos hacia el host lo que debe ser accesible desde fuera.
    \end{itemize}
    \begin{alertblock}{Idea clave}
      Dentro de la red Docker, el nombre del servicio funciona como DNS interno.
    \end{alertblock}
    \column{0.38\textwidth}
    \centering
    {\color{OTIred}\fontsize{46}{52}\selectfont\faSitemap}\\[0.15cm]
    {\large\textbf{DNS interno}}
  \end{columns}
\end{frame}

\begin{frame}{Red interna vs puerto publicado}
  \begin{center}
  \begin{tikzpicture}[
    box/.style={draw, rounded corners=6pt, minimum width=3.2cm, minimum height=1.25cm, align=center, text width=3.0cm, font=\small},
    arr/.style={-{Stealth[length=3mm]}, thick, dockerblue}
  ]
    \node[box, fill=UUlightgray, draw=UUdarkgray] (host) {\faLaptop\\\textbf{Host}\\localhost:5000};
    \node[box, fill=blue!10, draw=blue!60, right=0.8cm of host] (web) {\faPython\\\textbf{web}\\5000 interno};
    \node[box, fill=green!12, draw=green!60!black, right=0.8cm of web] (db) {\faDatabase\\\textbf{db}\\5432 interno};
    \draw[arr] (host) -- node[above, font=\scriptsize] {ports} (web);
    \draw[arr] (web) -- node[above, font=\scriptsize] {db:5432} (db);
  \end{tikzpicture}
  \end{center}
  \vspace{0.15cm}
  \begin{block}{Regla práctica}
    Publica la app web si necesitas navegador. No publiques la base de datos si solo la consume otro contenedor.
  \end{block}
\end{frame}

\begin{frame}[fragile]{Definir una red explícita en Compose}
\begin{lstlisting}[style=terminal, basicstyle=\ttfamily\tiny]
services:
  web:
    build: .
    networks:
      - backend

  db:
    image: postgres:16
    networks:
      - backend

networks:
  backend:
\end{lstlisting}
  \vspace{-0.05cm}
  \begin{alertblock}{Cuándo hacerlo}
    Define redes explícitas para separar tráfico o documentar mejor la arquitectura.
  \end{alertblock}
\end{frame}

\begin{frame}[fragile]{Comandos útiles para redes}
  \footnotesize
  \begin{tabular}{p{0.34\textwidth} p{0.55\textwidth}}
    \toprule
    Comando & Uso \\
    \midrule
    \texttt{docker network ls} & Lista redes existentes \\
    \texttt{docker network inspect <red>} & Muestra contenedores conectados \\
    \texttt{docker compose ps} & Relaciona servicios y puertos publicados \\
    \texttt{docker compose exec web sh} & Permite probar conectividad desde el servicio \\
    \bottomrule
  \end{tabular}
  \vspace{0.2cm}
  \begin{block}{Prueba rápida}
    Desde \texttt{web}, el host de PostgreSQL debe ser \texttt{db}, no \texttt{localhost}.
  \end{block}
\end{frame}

\MainHeaderSlide{Volúmenes y\\persistencia}{Mantener datos aunque el contenedor se elimine}
\section{S4 · Volúmenes y Persistencia}

\begin{frame}{Por qué se pierden datos}
  \begin{columns}[c]
    \column{0.55\textwidth}
    \footnotesize
    \begin{itemize}
      \item Un contenedor es reemplazable: puede eliminarse y recrearse.
      \item Los datos escritos dentro del filesystem del contenedor no son una estrategia durable.
      \item Las bases de datos necesitan almacenamiento externo al ciclo de vida del contenedor.
      \item Docker ofrece volúmenes y bind mounts para resolverlo.
    \end{itemize}
    \vspace{-0.05cm}
    \begin{alertblock}{Regla mental}
      Contenedor reemplazable; datos persistentes.
    \end{alertblock}
    \column{0.39\textwidth}
    \centering
    {\color{OTIred}\fontsize{46}{52}\selectfont\faDatabase}\\[0.15cm]
    {\large\textbf{Datos fuera del contenedor}}
  \end{columns}
\end{frame}

\begin{frame}{Volumen nombrado vs bind mount}
  \footnotesize
  \begin{tabular}{p{0.22\textwidth} p{0.35\textwidth} p{0.33\textwidth}}
    \toprule
    Tipo & Qué es & Cuándo usarlo \\
    \midrule
    Volumen nombrado & Almacenamiento administrado por Docker & Datos de bases de datos \\
    Bind mount & Carpeta local montada en el contenedor & Desarrollo y edición de código \\
    Sin volumen & Datos dentro del contenedor & Solo pruebas descartables \\
    \bottomrule
  \end{tabular}
  \vspace{0.2cm}
  \begin{block}{Analogía}
    El contenedor es una laptop prestada; el volumen es el disco externo donde guardas lo importante.
  \end{block}
\end{frame}

\begin{frame}[fragile]{Volumen nombrado para PostgreSQL}
\begin{lstlisting}[style=terminal, basicstyle=\ttfamily\scriptsize]
services:
  db:
    image: postgres:16
    volumes:
      - postgres_data:/var/lib/postgresql/data

volumes:
  postgres_data:
\end{lstlisting}
  \begin{alertblock}{Clave}
    Si eliminas el contenedor, el volumen puede seguir existiendo y conservar los datos.
  \end{alertblock}
\end{frame}

\begin{frame}[fragile]{Bind mount para desarrollo}
\begin{lstlisting}[style=terminal, basicstyle=\ttfamily\scriptsize]
services:
  web:
    build: .
    volumes:
      - ./app:/app
    ports:
      - "5000:5000"
\end{lstlisting}
  \begin{block}{Uso típico}
    Montar código local dentro del contenedor para desarrollo. Para producción, normalmente se prefiere copiar el código en la imagen.
  \end{block}
\end{frame}

\begin{frame}[fragile]{Comandos útiles para volúmenes}
  \footnotesize
  \begin{tabular}{p{0.34\textwidth} p{0.55\textwidth}}
    \toprule
    Comando & Uso \\
    \midrule
    \texttt{docker volume ls} & Lista volúmenes \\
    \texttt{docker volume inspect <volumen>} & Muestra detalles del volumen \\
    \texttt{docker compose down} & Borra contenedores y red, conserva volúmenes \\
    \texttt{docker compose down -v} & Borra también los volúmenes del proyecto \\
    \bottomrule
  \end{tabular}
  \vspace{0.2cm}
  \begin{alertblock}{Cuidado}
    \texttt{down -v} elimina datos persistentes del stack. Úsalo solo si realmente quieres reiniciar desde cero.
  \end{alertblock}
\end{frame}

\MainHeaderSlide{Servicios internos}{Aislar bases de datos y componentes privados}
\section{S4 · Aislamiento de Servicios}

\begin{frame}{Qué servicios deben ser internos}
  \begin{columns}[c]
    \column{0.52\textwidth}
    \small
    \begin{itemize}
      \item La app web suele publicar un puerto para el navegador.
      \item La base de datos normalmente solo debe recibir tráfico desde la app.
      \item Redis, workers y colas suelen ser servicios internos.
      \item Menos puertos publicados significa menor superficie de exposición.
    \end{itemize}
    \begin{alertblock}{Regla práctica}
      Si el usuario no necesita entrar directamente, no publiques el puerto.
    \end{alertblock}
    \column{0.42\textwidth}
    \centering
    \begin{tikzpicture}[
      box/.style={draw, rounded corners=6pt, minimum width=3cm, minimum height=1cm, align=center, font=\small},
      arr/.style={-{Stealth[length=3mm]}, thick, OTIred}
    ]
      \node[box, fill=UUlightgray, draw=UUdarkgray] (host) {Host};
      \node[box, fill=blue!10, draw=blue!60, below=0.45cm of host] (web) {web publicado};
      \node[box, fill=green!12, draw=green!60!black, below=0.45cm of web] (db) {db interno};
      \draw[arr] (host) -- node[right, font=\scriptsize] {5000} (web);
      \draw[arr] (web) -- node[right, font=\scriptsize] {5432} (db);
    \end{tikzpicture}
  \end{columns}
\end{frame}

\begin{frame}[fragile]{Base de datos interna en Compose}
\begin{lstlisting}[style=terminal, basicstyle=\ttfamily\scriptsize]
services:
  web:
    build: .
    ports:
      - "5000:5000"
    depends_on:
      - db

  db:
    image: postgres:16
    # Sin ports: no se expone al host
\end{lstlisting}
  \begin{block}{Lectura}
    \texttt{web} está publicado para el navegador. \texttt{db} solo está disponible dentro de la red Docker del proyecto.
  \end{block}
\end{frame}

\MainHeaderSlide{Healthchecks y backup}{Validar, respaldar y restaurar datos}
\section{S4 · Healthchecks y Backup}

\begin{frame}{Por qué usar healthchecks}
  \begin{columns}[c]
    \column{0.55\textwidth}
    \small
    \begin{itemize}
      \item Un contenedor puede estar encendido pero no listo.
      \item PostgreSQL puede tardar segundos en aceptar conexiones.
      \item Un healthcheck permite declarar cómo verificar el estado real.
      \item Mejora la observabilidad durante laboratorios y despliegues.
    \end{itemize}
    \begin{alertblock}{Idea clave}
      \texttt{running} no siempre significa \texttt{healthy}.
    \end{alertblock}
    \column{0.39\textwidth}
    \centering
    {\color{UUgreen}\fontsize{46}{52}\selectfont\faHeartbeat}\\[0.15cm]
    {\large\textbf{Listo para operar}}
  \end{columns}
\end{frame}

\begin{frame}[fragile]{Healthcheck para PostgreSQL}
\begin{lstlisting}[style=terminal, basicstyle=\ttfamily\scriptsize]
services:
  db:
    image: postgres:16
    healthcheck:
      test: ["CMD-SHELL", "pg_isready -U appuser -d appdb"]
      interval: 10s
      timeout: 5s
      retries: 5
\end{lstlisting}
  \begin{block}{Lectura}
    Docker ejecuta \texttt{pg\_isready}. Si responde correctamente, marca el servicio como saludable.
  \end{block}
\end{frame}

\begin{frame}[fragile]{Backup básico de PostgreSQL}
\begin{lstlisting}[style=terminal, basicstyle=\ttfamily\scriptsize]
# Crear carpeta local para backups
mkdir backups

# Exportar datos desde el servicio db
docker compose exec -T db pg_dump \
  -U appuser appdb > backups/appdb.sql
\end{lstlisting}
  \begin{alertblock}{Regla práctica}
    Un volumen conserva datos localmente, pero backup significa tener una copia exportable fuera del contenedor.
  \end{alertblock}
\end{frame}

\begin{frame}[fragile]{Restaurar información}
\begin{lstlisting}[style=terminal, basicstyle=\ttfamily\scriptsize]
# Restaurar desde el archivo SQL
docker compose exec -T db psql \
  -U appuser appdb < backups/appdb.sql
\end{lstlisting}
  \begin{block}{Antes de restaurar}
    Verifica que el servicio \texttt{db} esté arriba, que el usuario sea correcto y que el archivo exista en tu máquina.
  \end{block}
\end{frame}

\MainHeaderSlide{Laboratorio\\datos y redes}{Aislar PostgreSQL, persistir datos y probar respaldo}
\section{S4 · Laboratorio Datos y Redes}

\begin{frame}{Flujo del laboratorio}
  \begin{center}
  \begin{tikzpicture}[
    labstep/.style={draw, rounded corners=6pt, minimum width=2.35cm, minimum height=1.1cm, align=center, text width=2.15cm, font=\scriptsize},
    arr/.style={-{Stealth[length=3mm]}, thick, OTIred}
  ]
    \node[labstep, fill=blue!10, draw=blue!60] (s1) {\faNetworkWired\\\textbf{Red}\\interna};
    \node[labstep, fill=green!12, draw=green!60!black, right=0.45cm of s1] (s2) {\faDatabase\\\textbf{Volumen}\\persistente};
    \node[labstep, fill=orange!15, draw=orange!70, right=0.45cm of s2] (s3) {\faHeartbeat\\\textbf{Healthcheck}\\db};
    \node[labstep, fill=OTIred!10, draw=OTIred!70, right=0.45cm of s3] (s4) {\faSave\\\textbf{Backup}\\SQL};
    \draw[arr] (s1) -- (s2);
    \draw[arr] (s2) -- (s3);
    \draw[arr] (s3) -- (s4);
  \end{tikzpicture}
  \end{center}
  \begin{block}{Meta}
    Mejorar el stack Flask + PostgreSQL para que la base de datos sea interna, persistente y respaldable.
  \end{block}
\end{frame}

\begin{frame}[fragile]{Comandos del laboratorio}
  \footnotesize
  \begin{columns}[T]
    \column{0.50\textwidth}
\begin{lstlisting}[style=terminal]
# 1. Levantar stack
docker compose up -d --build

# 2. Revisar puertos
docker compose ps

# 3. Inspeccionar red
docker network ls
docker network inspect <red>
\end{lstlisting}
    \column{0.46\textwidth}
\begin{lstlisting}[style=terminal]
# 4. Ver volumenes
docker volume ls

# 5. Backup
mkdir backups
docker compose exec -T db pg_dump \
  -U appuser appdb > backups/appdb.sql
\end{lstlisting}
  \end{columns}
\end{frame}

\begin{frame}{Errores frecuentes en datos y redes}
  \scriptsize
  \begin{tabular}{p{0.28\textwidth} p{0.32\textwidth} p{0.32\textwidth}}
    \toprule
    Error & Causa probable & Solución rápida \\
    \midrule
    App no conecta a BD & Host incorrecto o red distinta & Usar \texttt{db} y misma red \\
    BD expuesta al host & Se agregó \texttt{ports} en \texttt{db} & Quitar \texttt{ports} si no es necesario \\
    Datos desaparecen & Sin volumen persistente & Agregar volumen nombrado \\
    Backup falla & Usuario o BD incorrectos & Revisar \texttt{DB\_USER} y \texttt{DB\_NAME} \\
    Servicio no está listo & Falta healthcheck o espera & Revisar logs y estado healthy \\
    \bottomrule
  \end{tabular}
\end{frame}

\begin{frame}{Checklist de aprendizaje — Sesión 4}
  \begin{itemize}
    \item Puedo explicar cómo Docker resuelve nombres de servicios en una red.
    \item Puedo diferenciar puerto publicado y comunicación interna.
    \item Puedo usar volúmenes nombrados para persistir datos.
    \item Puedo usar bind mounts para desarrollo local.
    \item Puedo mantener PostgreSQL como servicio interno.
    \item Puedo crear un backup y restaurarlo con comandos básicos.
  \end{itemize}
\end{frame}

\begin{frame}{Resumen de la Sesión 4}
  \footnotesize
  \begin{enumerate}
    \item Los servicios se comunican por nombre dentro de redes Docker.
    \item Los puertos publicados son para acceso desde el host, no para comunicación interna.
    \item Los volúmenes nombrados conservan datos aunque se recree el contenedor.
    \item Los bind mounts son útiles para desarrollo y edición local.
    \item Los healthchecks verifican disponibilidad real, no solo ejecución.
    \item Backup y restauración convierten datos persistentes en datos recuperables.
  \end{enumerate}
  \vspace{0.08cm}
  \begin{alertblock}{Próxima sesión}
    \textbf{Docker en producción:} Nginx, multi-stage builds, credenciales seguras, logging y monitoreo.
  \end{alertblock}
\end{frame}

% ============================================================
% CIERRE GENERAL
% ============================================================

\begin{frame}{Resumen del programa (hasta ahora)}
  \tiny
  \begin{enumerate}
    \item Docker resuelve el problema de reproducibilidad entre entornos.
    \item Contenerizar es empaquetar app + dependencias para ejecución uniforme.
    \item Una VM virtualiza una máquina completa; un contenedor comparte el kernel del host.
    \item La instalación de Docker se guía por la documentación oficial, no de memoria.
    \item Imagen = plantilla inmutable; contenedor = instancia en ejecución.
    \item Comandos esenciales: run, ps, stop, rm, images, build.
    \item Un Dockerfile es la receta profesional para construir imágenes propias.
    \item Las capas son inmutables: el orden correcto acelera los builds con caché.
    \item Slim y Alpine reducen drásticamente el tamaño de las imágenes.
    \item Docker Hub permite compartir imágenes con tagging y versionado.
    \item Docker Compose permite levantar stacks multi-contenedor reproducibles.
    \item Flask y PostgreSQL pueden comunicarse usando nombres de servicio.
    \item Redes Docker separan comunicación interna de puertos publicados.
    \item Volúmenes, healthchecks y backups hacen el stack más durable.
  \end{enumerate}
  \vspace{-0.02cm}
  \begin{alertblock}{Próxima sesión}
    \textbf{Docker en producción:} reverse proxy, credenciales, logging y monitoreo.
  \end{alertblock}
\end{frame}

\end{document}
